% concatenated from f2059.tex
\documentclass[a4paper]{amsart}
%\documentclass[a4paper]{amsart}
\usepackage{amssymb,latexsym}
\usepackage[utf8x]{inputenc}
\usepackage{amsmath,amsthm}
\usepackage{amsfonts,mathrsfs}
\usepackage{hyperref}
\usepackage{cleveref}
\usepackage{enumerate,units}
\usepackage[all]{xy}
\usepackage{graphicx,url}

%\usepackage{showlabels}
%\usepackage[left, pageswise, mathlines]{lineno}
%   \linenumbers
   
\theoremstyle{plain}
\newtheorem{theorem}{Theorem}[section]
\newtheorem{maintheorem}[theorem]{Main Theorem}
\newtheorem{corollary}[theorem]{Corollary}
\newtheorem{lemma}[theorem]{Lemma}
\newtheorem{proposition}[theorem]{Proposition}
\newtheorem{conclusion}[theorem]{Conclusion}
\newtheorem*{proposition*}{Proposition}
\newtheorem{remarkcnt}[theorem]{Remark}
\newtheorem{fact}[theorem]{Fact}
\newtheorem{claim}{Claim}[theorem]
\newtheorem*{subclaim}{Subclaim}
\newtheorem*{theorem*}{Theorem}
\newtheorem*{context*}{Assumption $\diamondsuit$}

\theoremstyle{definition}
\newtheorem{definition}[theorem]{Definition}
\newtheorem{assumption}[theorem]{Assumption}
\newtheorem{example}[theorem]{Example}

\theoremstyle{remark}
\newtheorem{remark}[theorem]{Remark}
\newtheorem*{notation}{Notation}
\newtheorem{question}[theorem]{\textbf{Question}}
\newtheorem{conjecture}[theorem]{\textbf{Conjecture}}


\numberwithin{equation}{section}
\newcommand{\forkindep}[1][]{%
  \mathrel{
    \mathop{
      \vcenter{
        \hbox{\oalign{\noalign{\kern-.3ex}\hfil$\vert$\hfil\cr
              \noalign{\kern-.7ex}
              $\smile$\cr\noalign{\kern-.3ex}}}
      }
    }\displaylimits_{#1}
  }
}

\newenvironment{claimproof}[1][\proofname]
  {%
    \proof[#1]%
      \renewcommand*\qedsymbol{$\square$ (claim)}%
  }
  {%
    \endproof%
  }

  
\newenvironment{subclaimproof}[1][\proofname]
  {%
    \proof[#1]%
      \renewcommand*\qedsymbol{‌​$\square$ (subclaim)}%
  }
  {%
    \endproof%
  }


   %% A "step" newenvironment
\newcounter{step}                   %defines a counter that will count the steps
\newenvironment{step}
    {\refstepcounter{step}
     \vspace{7pt}\par\noindent
     \textbf{Step \thestep.}\hspace{0.3mm}}%
    {\hfill $\clubsuit$             %puts a club-symbol at the end of the step
     \vspace{7pt}\par}

\makeatletter
\providecommand*{\cupdot}{%
  \mathbin{%
    \mathpalette\@cupdot{}%
  }%
}
\newcommand*{\@cupdot}[2]{%
  \ooalign{%
    $\m@th#1\cup$\cr
    \hidewidth$\m@th#1\cdot$\hidewidth
  }%
}
\makeatother

\newcommand{\reg}{\mathrm{reg}}
\newcommand{\dom}{\mathrm{Dom}}
\newcommand{\U}{\underline}
\newcommand{\Ba}{\textbf{a}}
\newcommand{\Bb}{\textbf{b}}
\newcommand{\Bc}{\textbf{c}}
\newcommand{\kap}{\varkappa}
\newcommand{\Sh}{\text{Sh}}
\newcommand{\LSh}{\text{LSh}}
\newcommand{\RSh}{\text{RSh}}
\newcommand{\Cyc}{\text{Cyc}}
\newcommand{\E}{\mathrel{E}}
\newcommand{\R}{\mathrel{R}}
\newcommand{\D}{\mathrel{D}}
\newcommand{\suc}[1]{\mathrm{Suc}(#1)}

\DeclareMathOperator{\acl}{acl}
\DeclareMathOperator{\dcl}{dcl} 
\DeclareMathOperator{\tp}{tp}
\DeclareMathOperator{\otp}{otp}
\DeclareMathOperator{\qftp}{qftp}
\DeclareMathOperator{\stp}{stp}
\DeclareMathOperator{\rk}{rk}
\DeclareMathOperator{\sgn}{sgn} 
\DeclareMathOperator{\st}{st}
\DeclareMathOperator{\Ur}{U}
\DeclareMathOperator{\Rg}{Range}
\DeclareMathOperator{\Img}{Im}
\DeclareMathOperator{\cf}{cf}
\DeclareMathOperator{\ch}{ch_T}
\DeclareMathOperator{\Sk}{Sk}


%************************************************************


%
%
%\newwrite\refs
%\openout\refs=\jobname.refs
%\makeatletter
%\renewcommand\@setref[3]{%
%        \ifx#1\relax
%                \write\refs{'#3' \thepage\space undefined}%
%                \protect \G@refundefinedtrue
%                \nfss@text{\reset@font\bfseries ??}%
%                \@latex@warning{Reference `#3' on page \thepage\space
%                                undefined}%
%        \else
%                \begingroup
%                \count@\expandafter\@secondoftwo#1\relax
%                \ifnum\c@page<\count@
%                        \write\refs{'#3' \thepage\space
%                                    \expandafter\@secondoftwo#1}%
%                \fi
%                \endgroup
%                \expandafter#2#1\null
%        \fi
%}
%\makeatother
  
\begin{document}
\title{Infinite Cliques in Simple and Stable Graphs}
%\author{Yatir Halevi}
%\email{yatirh@gmail.com}
%\begin{abstract}

%\end{abstract}
\author{Yatir Halevi}
\address{Department of Mathematics\\ University of Haifa\\ 199 Abba Khoushy Avenue \\ Haifa \\Israel}
 \email{ybenarih@campus.haifa.ac.il}

\author{Itay Kaplan}
\address{Einstein Institute of Mathematics, Hebrew University of Jerusalem, 91904, Jerusalem
Israel.}
\email{kaplan@math.huji.ac.il}

\author{Saharon Shelah}
\address{Einstein Institute of Mathematics, Hebrew University of Jerusalem, 91904, Jerusalem
Israel.}
\email{shelah@math.huji.ac.il}

\thanks{The first author would like to thanks the Israel Science Foundation for its support of this research (grant No. 555/21 and 290/19). The second author would like to thank the Israel Science Foundation for their support of this research (grant no. 1254/18). The third author would like to thank the Israel Science Foundation grants no. 1838/19 and  2320/23. This is Paper no. 1211 in the third author's  publication list.}

\keywords{chromatic number; stable graphs; simple graphs; infinite cliques; Taylor's Conjecture}
\subjclass[2020]{03C45; 05C15; 03C50}

\begin{abstract}
  Suppose that $G$ is a graph of cardinality $\mu^+$ with  chromatic number $\chi(G)\geq \mu^+$. One possible reason that this could happen is if $G$ contains a clique of size $\mu^+$. We prove that this is indeed the case when the edge relation is stable. When $G$ is a random graph (which is simple but not stable), this is not true. But still if in general the complete theory of $G$ is simple, $G$ must contain finite cliques of unbounded sizes.
\end{abstract}
  

\maketitle

\section{Introduction}
The chromatic number $\chi(G)$ of a graph $G$ is the minimal cardinal $\kappa$ for which there exists a vertex coloring with $\kappa$ colors such that connected vertices get different colors. 

Research around graphs having an uncountable chromatic number has a long history, see e.g., \cite[Section 3]{komjath}. This topic is set-theoretic in nature, with many results being independent of the axioms of set theory (ZFC). In \cite{1196,1211} we studied a specific conjecture (Taylor's strong conjecture) in the context of \emph{stable} graphs (in ZFC): a graph whose first order theory is stable (a model-theoretic notion of tameness, see Section \ref{sec:Preliminaries} for all the definitions). It turns out that a very close relative of this conjecture holds for stable graphs (although it does not hold in general).

More specifically, we showed that if a stable graph has chromatic number $>\beth_2(\aleph_0)$ then this implies the presence of all the finite subgraphs of a shift graph $\Sh_n(\omega)$ for some $0<n < \omega$, where for a cardinal $\kappa$, the shift graph $\Sh_n(\kappa)$ is the graph whose vertices are increasing $n$-tuples of ordinals in $\kappa$, and two such tuples $s,t$ are connected if for every $1\leq i\leq n-1$, $s(i)=t(i-1)$ or vice-versa (see Example \ref{E:shift}). In turn, this implies that the chromatic numbers of elementary extensions of said graph are unbounded. An important example for this paper is the case $n=1$: $\Sh_1(\kappa)$ is the complete graph on $\kappa$. For more on this result, see \cite{1196,1211}. 

In this paper, we take the first step towards identifying the $n$ in the previous paragraph, by considering not only the chromatic number of the graph, but its \emph{cardinality} as well, as we now explain.

Bounds for the chromatic number of the shift graph was computed by Erd\"os and Hajnal \cite{EH}: assuming the generalized continuum hypothesis (\emph{GCH}), 
\[\chi(\Sh_{n+1}(\kappa^{+n}))=\kappa\]
for all $n<\omega$, see Fact \ref{F:Sh-high chrom}. In particular (and trivially) $\chi(\Sh_1(\kappa))= \kappa$. Thus, it makes sense to ask the following question:

\begin{question} \label{Q:The question on Ch}
  Suppose that $G$ is a stable graph and for simplicity also GCH. Assume that for every cardinal $\kappa$, there is some $G' \equiv G$ (i.e., $\mathrm{Th}(G)=\mathrm{Th}(G')$) of cardinality $|G'| \leq \kappa^{+n}$ satisfying that $\chi(G') \geq \kappa^+$, then is it true that for some $m \leq n$, $G$ contains all finite subgraphs of $\Sh_{m}(\omega)$?
\end{question}

\begin{remark}
  Proposition \ref{P: stable no arbt large finite cliques} and Remark \ref{R:even stable theory} explain why we restricted ourselves to successor cardinals. 
\end{remark}

In this paper we deal with the case $n=1$, and we manage to give a satisfying solution to this case (and more) assuming that the theory of $G$ is simple or that the edge relation is stable (both weaker assumptions than stability of the theory). 
The following sums up the main results of this paper: 

\begin{maintheorem}[Propositions \ref{P:cannot embed infinite complete graph} and \ref{P:stable}]
  Let $G=(V,E)$ be a graph and $T=\mathrm{Th}(G,E)$ be its first order theory. Assume that $|G|=\mu^+$ and $\chi(G)\geq \mu^+$ for some infinite cardinal $\mu$.
  
  \begin{enumerate}
  \item Assuming $T$ is a simple theory then $G$ contains cliques of any finite size. 
  \item Assuming the edge relation $E$ is stable then $G$ contains an infinite clique of cardinality $\mu^+$.
  \end{enumerate}
  \end{maintheorem}

Note that the conclusion of item (1) (together with L\"owenheim-Skolem and compactness) implies the existence of $G'\succ G$ of cardinality $\mu^+$ that contains an infinite clique of cardinality $\mu^+$. The conclusion of item (2) is stronger: we can find such a clique already in $G$ itself.

\subsection{Organization of the paper} In Section \ref{sec:Preliminaries} we go over the relevant basic definitions in model theory and graph theory. Section \ref{sec:infinite cliques} contains the proof of the main results. Section \ref{S:chromatic} reframes the main results in terms of the function hinted on in Question \ref{Q:The question on Ch}. Finally, in the Appendix \ref{A:Hajnal-Komjath example} we analyze an example of Hajnal and Komj\'{a}th that was given by them to refute Taylor's strong conjecture and show that the theory of this graph is not stable (in fact, we show more: that it is not simple and has IP).
% A particular example  $\Sh_n(\kappa)$ is an infinite complete graph it is interesting to investigate which ``tame'' graphs contain all the finite complete graphs. 

% Consider a graph $G=(V,E)$ for which $\chi(G)\geq \mu$ for some infinite cardinal $\mu$ (so $|G|\geq \mu$ as well) and assume that it contains arbitrary large finite cliques. An easy compactness and L\"owenheim-Skolem argument shows the existence of some $G'\succ G$ with $|G'|=|G|$ that contains an infinite clique of cardinality $\mu$. This cannot be hoped to exist for any cardinal $\mu$, there exists a graph $G$ of cardinality $\mu$ (with $\mu$ not a successor) with $\chi(G)\geq \mu$ that contains no arbitrary large finite cliques (Proposition \ref{P: stable no arbt large finite cliques}), so we restrict to successor cardinals. 


% On the other hand, it is not hard to see that some model theoretic ``tameness'' assumption is needed to guarantee the existence of arbitrary large finite cliques.  
% The generic triangle-free graph is the Fra\"iss\'e limit of all triangle-free finite graphs. Its theory $T$ has TP$_2$, SOP$_3$ and NSOP$_4$. 
% Now, let $G$ be a triangle-free graph with $|G|=\mu$ and $\chi(G)=\mu$, see \cite{ErRa-triangle-free} for an example of such a graph. So $G$ can be embedded inside a sufficiently saturated model $\mathbb{H}_2\vDash T$. By L\"owenheim-Skolem, we can find a model $H_2\vDash T$ of cardinality $\mu$ which contains $G$. As a result $\chi(H_2)=\mu$ and is triangle-free.

% Furthermore, assuming CH, the Hanjanl-Komj\'ath example from \cite{HK}, is an example of such a graph (of cardinality $\aleph_1$), which does not contain all the finite subgraphs of any shift graph. In Proposition \ref{P:Hajnal-Komjath Example} we show that their example has IP.

% The question is:
% \begin{question}
% Let $\mu$ be an infinite cardinal and $G$ a graph with $\chi(G)\geq \mu^+$ and $|G|=\mu^+$. Which ``tameness'' conditions guarantee that $G$ contains arbitrary finite cliques?\footnote{In \cite{KoSh}, Komj\'ath and the third author  characterized the finite subgraphs of such graphs.}
% \end{question}

% Our main results of the paper answer this question for \emph{stable}  graphs and \emph{simple} graphs.

% Stability theory, which is the study of stable theories and originated in the works of the third author in the 60s and 70s, is one of the most influential and important subjects in modern model theory. Examples of stable theories include abelian groups, modules, algebraically closed fields, graph theoretic trees, or more generally superflat graphs \cite{PZ}. Stability also had an impact in combinatorics, e.g. \cite{MS} and \cite{CPT} to name a few. 

% Simplicity, which also originated in the works of the third author, encompasses examples like pseudo-finite fields, the model companion of fields with automorphisms (i.e. ACFA) and the random graph (i.e. the Rado graph). Every stable structure is a simple structure.



% The main result of this paper is the following.



%Taylor's conjecture for stable graphs was studied in \cite{1196,1211}. It was refuted in general by Hajnal and Komjath in \cite{HK}. In Appendix we investigate their example and show that it has the \emph{independence property}. On the other hand, in \cite{1196}, we showed that it holds for $\omega$-stable graphs. In it interesting to find an example which is stable.
%

\section{Preliminaries}\label{sec:Preliminaries}

We use small latin letters $a,b,c$ for tuples and capital letters $A,B,C$ for sets. We also employ the standard model theoretic abuse of notation and write $a\in A$ even for tuples when the length of the tuple is immaterial or understood from context.

\subsection{Stability and Simplicity}
We use fairly standard model theoretic terminology and notation, see for example \cite{TZ}. We gather some of the needed notions. For stability, the reader may also consult with \cite{classification}.

We denote by $\tp(a/A)$ the complete type of $a$ over $A$. 
%Let $(I,<)$ be a linearly ordered set. A sequence $\langle a_i: i\in I\rangle$ inside a first order structure is \emph{indiscernible} if for any $i_1<\dots<i_k$ and $j_1<\dots<j_k$ in $I$,
%\[\tp(a_{i_1},\dots,a_{i_k})=\tp(a_{j_1},\dots, a_{j_k}).\]
A structure $M$ is \emph{$\kappa$-saturated}, for a cardinal $\kappa$, if any type $p$ over $A$ with $|A|< \kappa$ is realized in $M$. The structure $M$ is \emph{saturated} if it is $|M|$-saturated.

The \emph{monster model} of a complete theory $T$, denoted here by $\mathbb{U}$, is a large saturated model containing all sets and models (as elementary substructures) we  encounter.\footnote{There are set theoretic issues in assuming that such a model exists, but these are overcome by standard techniques from set theory that ensure the generalized continuum hypothesis from some point on while fixing a fragment of the universe, see \cite{HaKa}. The reader can just accept this or alternatively assume that $\mathbb{U}$ is merely $\kappa$-saturated and $\kappa$-strongly homogeneous for large enough $\kappa$.}  All subsets and models will be \emph{small}, i.e. of cardinality $<|\mathbb{U}|$.

Given a first order theory $T$, a formula $\varphi(x,y)$ is $\emph{stable}$ if we cannot find elements $\langle a_i\in \mathbb{U}: i<\omega\rangle$ such that $\mathbb{U}\models \varphi(a_i,a_j)\iff i<j$. 

For any formula $\varphi(x,y)$ we set $\varphi(y,x)^{\text{opp}}=\varphi(x,y)$: it is the same formula but the roles of the variables replaced. The following is folklore.
\begin{fact}\label{F:def of types}
Work in a complete first order theory $T$ with infinite models. Let $\varphi(x,y)$ be a stable formula. There is a formula $\psi(y,z)$ such that  any $\varphi$-type $p$ over a model $M$ is definable by an instance of $\psi$ over $M$, i.e. for any such $\varphi$-type $p$ there is an element $c\in M$ such that $\varphi(x,b)\in p\iff M\vDash  \psi(b,c)$. Moreover, $\psi(y,z)$ can be chosen such that for any $c\in M$, $\psi(y,c)$ is equivalent to a boolean combination of instances of $\varphi^{\text{opp}}$. 
\end{fact}
\begin{proof}
By \cite[Theorem II.2.12]{classification}, there is some $\psi(y,z)$ such that for any $\varphi$-type $p$ over any set $A$ ($|A|\geq 2$), there is some $c_p\in A$ such that $\psi(y,c_p)$ defines $p$. By \cite[Lemma 2.2(i)]{Pillay}, if $p\in S_{\varphi}(M)$, where $M$ is a model, then $p$ is definable by a boolean combination of instances of $\varphi^{\text{opp}}$. But then as $M$ is a model $\psi(y,c_p)$ is equivalent to such a boolean combination.
\end{proof}

A theory $T$ is \emph{stable} if all formulas are stable.

Next we define simplicity; we give an equivalent definition, using the notion of dividing for types, see \cite[Proposition 7.2.5]{TZ}. Given a first order theory $T$ with a monster model $\mathbb{U}$, a formula $\varphi(x,b)$ with $b\in \mathbb{U}$ \emph{divides} over $A$ if there is a sequence of realizations $\langle b_i\in \mathbb{U}:i<\omega\rangle$ of $\tp(b/A)$ such that $\{\varphi(x,b_i):i<\omega\}$ is $k$-inconsistent for some $k<\omega$ (every set of size $k$ is inconsistent). A complete type $p$ over $B$ divides over $A$ if it contains some formula which divides over $A$.

The theory $T$ is \emph{simple} if for every complete type $p$ over $B$ there is some $A\subseteq B$ with $|A|\leq |T|$ such that $p$ does not divide over $A$. Every stable theory is simple.


The main tool we will use from simplicity theory is forking calculus. Non-forking independence is 3-place relation on sets (or tuples) denoted by $\forkindep$. We will not go over all the properties that non-forking independence enjoys in simple theories; see \cite[Chapter 7]{TZ} for more information.

\subsection{Graph theory}
Here we gather some facts on graphs and the chromatic number of graphs (see also \cite{1196}).

By a \emph{graph} we mean a pair $G=(V,E)$ where $E\subseteq V^2$ is symmetric and irreflexive. A \emph{graph homomorphism} between $G_1=(V_1,E_1)$ and $G_2=(V_2,E_2)$ is a map $f:V_1\to V_2$ such that $f(e)\in E_2$ for every $e\in E_1$. If $f$ is injective we will say that $f$ embeds $G_1$ into $G_2$ a subgraph. If in addition we require that $f(e)\in E_2$ if and only if $e\in E_1$ we will say that $f$ embeds $G_1$ into $G_2$ as an induced subgraph.


\begin{definition}
Let $G=(V,E)$ be a graph.
\begin{enumerate}
\item For a cardinal $\kappa$, a \emph{vertex coloring} (or just coloring) of cardinality $\kappa$ is a function $c:V\to \kappa$ such that $x\E y$ implies $c(x)\neq c(y)$ for all $x,y\in V$.
\item The \emph{chromatic number} $\chi(G)$ is the minimal cardinality of a vertex coloring of $G$.
\end{enumerate}
\end{definition}

The following is easy and well known.

\begin{fact}\label{F:basic-prop-chi}
Let $G=(V,\E)$ be a graph.
If $V=\bigcup_{i\in I} V_i$ then $\chi(G)\leq \sum_{i\in I} \chi(V_i, E\restriction V_i)$.
\end{fact}
\begin{proof}
Let $c_i:V_i\to \mu_i$ be a coloring of $(V_i,E\restriction V_i)$. Define a coloring $c:V\to \bigcup \{\mu_i\times\{i\}:i\in I\}$ by choosing for any $v\in V$ an $i_v\in I$ such that $v\in V_{i_v}$ and setting $c(v)=(c_{i_v}(v),i_v)$.
\end{proof}

\begin{example} \label{E:shift}
For any finite number $r\geq 1$ and any linearly ordered set $(A,<)$, let $\Sh_r(A)$, (\emph{the shift graph on $A$}) be the following graph: its set of vertices is the set of strictly increasing $r$-tuples from $A$, $\langle s_0<\dots<s_{r-1}\rangle $, and we put an edge between $s$ and $t$ if for every $1\leq i\leq r-1$, $s_i=t_{i-1}$, or vice-versa. It is an easy exercise to show that $\Sh_r(A)$ is a connected graph. If $r=1$ this gives $K_{A}$, the complete graph on $A$.
\end{example}


\begin{example}[Symmetric Shift Graph]
Let $r\geq 1$ be any natural number and $A$ any set. The \emph{symmetric shift graph} $\Sh_r^{sym}(A)$ is defined similarly as the shift graph but with set of vertices the set of distinct $r$-tuples.
Note that $\Sh_r(A)$ is an induced subgraph of $\Sh_r^{sym}(A)$ (and that for $r=1$ they are both the complete graph on $A$).

Since for any infinite set $A$, $\Sh_r^{sym}(A)$ is definable in $(A,=)$, it is stable. Moreover, for any two infinite sets $A$ and $B$, $\Sh_r^{sym}(A)\equiv \Sh_r^{sym}(B)$. Since every infinite set $A$ is saturated, $\Sh_r^{sym}(A)$ is saturated.
\end{example}


\begin{fact}\cite[Fact 2.6]{1196}\cite[Proof of Theorem 2]{EH-shift}\label{F:Sh-high chrom}
Let $1\leq r<\omega$ be a natural number and $\mu$ be an infinite cardinal,  \[\chi\left(\Sh_r^{sym}(\beth_{r-1}\left(\mu\right))\right)\leq\mu\] and \[\chi\left(\Sh_r(\beth_{r-1}\left(\mu\right)^{+})\right)\geq\mu^{+}.\]
\end{fact}

\begin{remark}
Note that it follows that $\chi(\Sh_{r}(\beth_{r-1}(\mu)))\leq \mu$.
\end{remark}


\begin{lemma}\label{L:sat of sym}
If $A$ is infinite then  $\Sh_r^{sym}(A)$ is a saturated model of $\mathrm{Th}(\Sh_r^{sym}(\omega))$ of cardinality $|A|$.
\end{lemma}
\begin{proof}
This follows easily from the fact that $(A,=)$ is saturated and that $\Sh_r^{sym}(A)$ is definable in $(A,=)$.
\end{proof}



We prove two easy results on the theory of the shift graphs.
\begin{lemma}\label{L:embed model of shift into symmetric}
Every model of  $T=\mathrm{Th}(\Sh_r(\omega))$ (of cardinality $\lambda$) can be embedded as an induced subgraph of $\Sh_r^{sym}(A)$ for some infinite set $A$ (of cardinality $\lambda$). 
\end{lemma}
\begin{proof}
Since $\Sh_r(\omega)$ is an induced subgraph of $\Sh_r^{sym}(\omega)$, the former satisfies the universal theory of $\Sh_r^{sym}(\omega)$. Thus every model $M$ of $T$ can be embedded as an induced subgraph of a model of $\mathrm{Th}(\Sh_r^{sym}(\omega))$. By Lemma \ref{L:sat of sym}, and the universality of saturated models, $M$ can be embedded as an induced subgraph of $\Sh_r^{sym}(A)$ for some infinite set $A$ of cardinality $|M|$.
\end{proof}


\begin{lemma}\label{L:properties of shift2}
For any cardinal $\mu$, the following holds for the graph $(\Sh_2(\mu),E)$:
\begin{enumerate}
\item Its complete theory is not stable.
\item Its graph relation is stable.
\item It is triangle-free.
\end{enumerate}
\end{lemma}
\begin{proof}
(1) For any $1<n<\omega$, let $X_n$ be the definable set $\{x\in \Sh_2(\mu): (x\E (0,n)) \wedge \neg(x\E (0,1))\}$, it is easily seen that $X_n=\{(n,k): k>n\}$. For any $0<m<\omega$, let $Y_m=\{x\in \Sh_2(\mu): (x\E (m,m+1))\wedge \neg(x\E (m-1,m+1))\}$, it is easily seen that $Y_m=\{(k,m): k<m\}$. 

For any pair of natural numbers $(n,m)$, with $n>1$ and $n<m$, let $\psi_{n,m}(x,y)$ be the formula (with parameters $n,m$): $\exists z (x\in X_n\wedge y\in Y_m\wedge (z\E x)\wedge (z\E y))$. 

For any $l$ with $l<m-n$ let $a_l=(n,n+l)\in X_n$, and for any $0<k<m-n$ let $b_k=(n+k,m)\in Y_m$. It is easily checked that $\psi_{n,m}(a_l,b_k) \iff l<k$. Since all the $\psi_{n,m}$'s are uniformly definable with parameters, by compactness we get that the theory of $\Sh_2(\mu)$ is not stable.

(2) One can either see directly that the graph relation is stable, or note that since $\Sh_2(\mu)$ is an induced subgraph of the stable graph $\Sh_2^{sym}(\mu)$ its edge relation must also be stable.

(3) Suppose that $(a,b)\E(c,d)\E(e,f)\E(a,b)$. Note that also $(c,d)\E(a,b)\E(e,f)\E(c,d)$. Hence, without loss of generality $b=c$. It follows easily that  $e=d$. So $a<b<d<f$. Now either $a=f$ or $b=e=d$, so either way we get a contradiction.
\end{proof}

\begin{remark}\label{R:Sh3}
\begin{enumerate} 
\item Since $\Sh_2(\mu)$ is definable in $(\mu,<)$, it has NIP (not the Independence Property (IP), see e.g., \cite{guidetonip}). 
\item One can also note that $\Sh_3^{sym}(\mu)$ is triangle-free.
\end{enumerate}
\end{remark}


Every graph is a (not necessarily induced) subgraph of a stable graph (e.g. a large enough complete graph). On the other hand, every shift graph $\Sh_n(\mu)$ is an induced subgraph of a stable graph (the symmetric shift graph). This raises the following:

\begin{question}
Is every graph with a stable edge relation an induced subgraph of a stable graph?
\end{question}

\section{Infinite Cliques}\label{sec:infinite cliques}
In this section we prove the main results of the paper. 

%We start with the following easy general observation which follows from compactness and L\"owenheim-Skolem; see also Section \ref{S:chromatic} for another viewpoint. 
%
%\begin{lemma}\label{L:background}
%Let $T$ be a complete theory of graphs. The following are equivalent:
%\begin{enumerate}
%\item $T$ proves the existence of arbitrary large finite cliques.
%\item There exists a model $M\vDash T$ which contains an infinite clique.
%\item For any $M\vDash T$ with $|M|=\mu$ there exists $N\succ M$, with $|N|=\mu$, which contains an infinite clique of cardinality $\mu$ (in particular $\chi(N)\geq \mu$). 
%\end{enumerate}
%\end{lemma}

\subsection{Simple Graphs}
We start with the following technical result.

\begin{lemma}\label{L:technical simple lemma}
Let $M$ be some structure, in a language $L$, and assume that $T=\mathrm{Th}(M)$ is simple. Let $M'$ be some expansion of $M$ to a language $L'\supseteq L$. Let $<$ be a linear order on the universe of $M$ and $\alpha,\beta\in M$ elements satisfying that
\begin{enumerate}
\item $(\dcl_{L'}(\alpha),<)$ and $(\dcl_{L'}(\beta),<)$ are well-orders and
\item  for any $\gamma\in \dcl_{L'}(\alpha)\cup \dcl_{L'}(\beta)$, $\gamma \forkindep[\dcl_{L'}(\gamma)\cap \{\varepsilon\in M:\varepsilon<\gamma\}]\{\varepsilon\in M:\varepsilon<\gamma\}$.
\end{enumerate}


Then  \[\dcl_{L'}(\alpha)\forkindep[\dcl_{L'}(\alpha)\cap \dcl_{L'}(\beta)]\dcl_{L'}(\beta).\]
\end{lemma}
\begin{proof}
Let $\alpha,\beta\in M$ be elements as in the statement, and let $\mathbf{a}^\alpha=\dcl_{L'}(\alpha)$ and $\mathbf{a}^\beta=\dcl_{L'}(\beta)$.
Set $\Omega=\mathbf{a}^\alpha\cup\mathbf{a}^\beta$ and for any $\gamma\in \Omega$ let $\Omega_{<\gamma}=\{\varepsilon \in \Omega: \varepsilon<\gamma\}$. Since $(\Omega,<)$ is a finite union of well ordered sets, it is also well ordered. 

\begin{claim}\label{C:induction}
For $\gamma\in \Omega$, if \[\mathbf{a}^\alpha\cap \Omega_{<\gamma} \forkindep[\mathbf{a}^\alpha\cap \mathbf{a}^\beta\cap \Omega_{<\gamma}]\mathbf{a}^\beta\cap\ \Omega_{<\gamma}\]
then
\[\mathbf{a}^\alpha\cap \Omega_{\leq \gamma} \forkindep[\mathbf{a}^\alpha\cap \mathbf{a}^\beta\cap \Omega_{\leq \gamma}]\mathbf{a}^\beta\cap\ \Omega_{\leq \gamma}.\]
\end{claim}
\begin{claimproof}
By symmetry, we deal with the case $\gamma\in \mathbf{a}^\alpha$. We first prove that
\begin{equation}\label{eq:indep}
\mathbf{a}^\alpha\cap \Omega_{\leq \gamma}\forkindep[\mathbf{a}^\alpha\cap \mathbf{a}^\beta\cap \Omega_{<\gamma}] \mathbf{a}^\beta\cap \Omega_{<\gamma}.
\end{equation}

By hypothesis (2), $\gamma\forkindep[\dcl_{L'}(\gamma)\cap \{\varepsilon:\varepsilon<\gamma\}]\{\varepsilon:\varepsilon<\gamma\}$.
Since \[\dcl_{L'}(\gamma)\cap\{\varepsilon:\varepsilon<\gamma\}\subseteq \mathbf{a}^\alpha\cap \Omega_{<\gamma}\subseteq \{\varepsilon:\varepsilon<\gamma\},\] we get that $\gamma\forkindep[\mathbf{a}^\alpha\cap \Omega_{<\gamma}]\{\varepsilon:\varepsilon<\gamma\}$, so $\gamma\forkindep[\mathbf{a}^\alpha\cap \Omega_{<\gamma}] \mathbf{a}^\beta\cap \Omega_{<\gamma}$ and \[\mathbf{a}^\alpha\cap\Omega_{\leq \gamma}\forkindep[\mathbf{a}^\alpha\cap \Omega_{<\gamma}] \mathbf{a}^\beta\cap \Omega_{<\gamma}.\] By assumption and transitivity, we conclude (\ref{eq:indep}).

If $\gamma\notin \mathbf{a}^\beta$ then we are done; so assume that $\gamma\in \mathbf{a}^\beta$ as well. In this case, by properties of forking we get that $\mathbf{a}^\alpha\cap  \Omega_{\leq \gamma}\forkindep[\mathbf{a}^\alpha\cap \mathbf{a}^\beta\cap \Omega_{\leq \gamma}]\mathbf{a}^\beta\cap\Omega_{\leq \gamma}$, which is what we wanted to prove.
\end{claimproof}
The proof now follows by induction: the successor step is Claim \ref{C:induction} and limit case follows since forking is witnessed by a formula.
%
%If $\gamma\notin \mathbf{a}^\alpha\cup \mathbf{a}^\beta$ then $\mathbf{a}^\alpha\cap (\gamma+1)=\mathbf{a}^\alpha\cap (\gamma\cup \{\gamma\})=\mathbf{a}^\alpha\cap \gamma$, and similarly for $\mathbf{a}^\beta$, so there is nothing to show. 
%
%Otherwise there are two options, either $\gamma\in \mathbf{a}^\alpha$ or $\gamma\in \mathbf{a}^\beta\setminus \mathbf{a}^\alpha$ or $\gamma\in \mathbf{a}^\alpha\cap \mathbf{a}^\beta$.
%
%Assume that $\gamma\in \mathbf{a}^\alpha\setminus \mathbf{a}^\beta$. By assumption we need to show that 
%Assume that $\gamma\in \mathbf{a}^\alpha\cap \mathbf{a}^\beta$. We want to prove that $(\mathbf{a}^\alpha\cap \gamma)\cup\{\gamma\}\forkindep[(\mathbf{a}^\alpha\cap \mathbf{a}^\beta\cap \gamma)] (\mathbf{a}^\beta\cap \gamma)$. As $\dcl_{L'}(\{\gamma\})\cap \gamma\subseteq \mathbf{a}^\alpha\cap \mathbf{a}^\beta\cap \gamma$ and $\{\gamma\}\forkindep[\dcl_{L'}(\{\gamma\})\cap\gamma]\gamma$ we get that $\{\gamma\}\forkindep[\mathbf{a}^\alpha\cap \mathbf{a}^\beta\cap \gamma] \gamma$ and thus $\{\gamma\}\forkindep[(\mathbf{a}^\alpha\cap \mathbf{a}^\beta\cap\gamma)\cup (\mathbf{a}^\alpha\cap \gamma)] (\mathbf{a}^\beta\cap \gamma)$, so  $\{\gamma\}\forkindep[\mathbf{a}^\alpha\cap \gamma] (\mathbf{a}^\beta\cap \gamma)$. Consequently, $(\mathbf{a}^\alpha\cap \gamma)\cup\{\gamma\}\forkindep[\mathbf{a}^\alpha\cap \gamma] (\mathbf{a}^\beta\cap \gamma)$. As above, it follows by the induction hypothesis and transitivity that $(\mathbf{a}^\alpha\cap \gamma)\cup\{\gamma\}\forkindep[(\mathbf{a}^\alpha\cap \mathbf{a}^\beta\cap\gamma)] (\mathbf{a}^\beta\cap \gamma)$.
\end{proof}



%
%%
%\begin{lemma}\label{L:technical simple lemma}
%Let $M=(\mu,\dots)$ be some structure, in a language $L$, on the ordinal $\mu$ and assume that $T=\mathrm{Th}(M)$ is simple. Let $M'$ be some expansion of $M$ to a language $L'\supseteq L$. Assume that for  all $\alpha\in M$, $\{\alpha\}\forkindep[\dcl_{L'}(\{\alpha\})\cap \alpha]\alpha$. 
%
%Then for any $\alpha,\beta \in M$, $\dcl_{L'}(\{\beta\})\forkindep[\dcl_{L'}(\{\alpha\})\cap \dcl_{L'}(\{\beta\})]\dcl_{L'}(\{\alpha\})$.
%\end{lemma}
%\begin{proof}
%Let $\alpha,\beta\in M$ and let $\mathbf{a}^\alpha=\dcl_{L'}(\{\alpha\})$ and $\mathbf{a}^\beta=\dcl_{L'}(\{\beta\})$.
%
%
%We will prove by induction on ordinals $\gamma<\mu$ that $\mathbf{a}^\alpha\cap \gamma \forkindep[\mathbf{a}^\alpha\cap \mathbf{a}^\beta\cap \gamma]\mathbf{a}^\beta\cap\ \gamma$; this follows since forking is witnessed by a formula.
%The limit ordinal case is proved similarly; so we assume it is know for $\gamma$ and prove it for $\gamma+1$.
%
%If $\gamma\notin \mathbf{a}^\alpha\cup \mathbf{a}^\beta$ then $\mathbf{a}^\alpha\cap (\gamma+1)=\mathbf{a}^\alpha\cap (\gamma\cup \{\gamma\})=\mathbf{a}^\alpha\cap \gamma$, and similarly for $\mathbf{a}^\beta$, so there is nothing to show. 
%
%First, assume that $\gamma\in \mathbf{a}^\alpha$. We prove that
%$(\mathbf{a}^\alpha\cap \gamma)\cup\{\gamma\}\forkindep[\mathbf{a}^\alpha\cap \mathbf{a}^\beta\cap \gamma] \mathbf{a}^\beta\cap \gamma$. By the hypothesis, $\{\gamma\}\forkindep[\dcl_{L'}(\{\gamma\})\cap \gamma]\gamma$ and thus $\{\gamma\}\forkindep[(\mathbf{a}^\alpha\cap \mathbf{a}^\beta\cap \gamma)\cup \{\dcl_{L'}(\{\gamma\})\cap \gamma\}] \gamma$. As $\gamma\in \mathbf{a}^\alpha$, $\dcl_{L'}(\{\gamma\})\cap \gamma \subseteq \mathbf{a}^\alpha\cap \gamma$, $\{\gamma\}\forkindep[(\mathbf{a}^\alpha\cap \mathbf{a}^\beta\cap \gamma)\cup (\mathbf{a}^\alpha\cap \gamma)] \gamma$ and $\{\gamma\}\forkindep[(\mathbf{a}^\alpha\cap \mathbf{a}^\beta\cap \gamma)\cup (\mathbf{a}^\alpha\cap \gamma)] \mathbf{a}^\beta\cap \gamma$, so $\{\gamma\}\forkindep[\mathbf{a}^\alpha\cap \gamma] \mathbf{a}^\beta\cap \gamma$. By the induction hypothesis and transitivity we conclude.
%
%If $\gamma\notin \mathbf{a}^\beta$ then we are done; so assume that $\gamma\in \mathbf{a}^\beta$ as well. In this case, by properties of forking in properties we get that $(\mathbf{a}^\alpha\cap \gamma)\cup \{\gamma\}\forkindep[(\mathbf{a}^\alpha\cap \mathbf{a}^\beta\cap \gamma)\cup \{\gamma\}](\mathbf{a}^\beta\cap \gamma)\cup\{\gamma\}$, which is what we wanted to prove.
%%
%%Otherwise there are two options, either $\gamma\in \mathbf{a}^\alpha$ or $\gamma\in \mathbf{a}^\beta\setminus \mathbf{a}^\alpha$ or $\gamma\in \mathbf{a}^\alpha\cap \mathbf{a}^\beta$.
%%
%%Assume that $\gamma\in \mathbf{a}^\alpha\setminus \mathbf{a}^\beta$. By assumption we need to show that 
%%Assume that $\gamma\in \mathbf{a}^\alpha\cap \mathbf{a}^\beta$. We want to prove that $(\mathbf{a}^\alpha\cap \gamma)\cup\{\gamma\}\forkindep[(\mathbf{a}^\alpha\cap \mathbf{a}^\beta\cap \gamma)] (\mathbf{a}^\beta\cap \gamma)$. As $\dcl_{L'}(\{\gamma\})\cap \gamma\subseteq \mathbf{a}^\alpha\cap \mathbf{a}^\beta\cap \gamma$ and $\{\gamma\}\forkindep[\dcl_{L'}(\{\gamma\})\cap\gamma]\gamma$ we get that $\{\gamma\}\forkindep[\mathbf{a}^\alpha\cap \mathbf{a}^\beta\cap \gamma] \gamma$ and thus $\{\gamma\}\forkindep[(\mathbf{a}^\alpha\cap \mathbf{a}^\beta\cap\gamma)\cup (\mathbf{a}^\alpha\cap \gamma)] (\mathbf{a}^\beta\cap \gamma)$, so  $\{\gamma\}\forkindep[\mathbf{a}^\alpha\cap \gamma] (\mathbf{a}^\beta\cap \gamma)$. Consequently, $(\mathbf{a}^\alpha\cap \gamma)\cup\{\gamma\}\forkindep[\mathbf{a}^\alpha\cap \gamma] (\mathbf{a}^\beta\cap \gamma)$. As above, it follows by the induction hypothesis and transitivity that $(\mathbf{a}^\alpha\cap \gamma)\cup\{\gamma\}\forkindep[(\mathbf{a}^\alpha\cap \mathbf{a}^\beta\cap\gamma)] (\mathbf{a}^\beta\cap \gamma)$.
%\end{proof}
We move to the main result of this section on simplicity.

\begin{proposition}\label{P:cannot embed infinite complete graph}
Let $T$ be a complete simple theory of graphs in the language of graphs $L=\{E\}$.

Let $\mu$ be an infinite cardinal and $G=(V,E)\vDash T$ with $|G|\leq 2^\mu$. If $\chi(G)\geq \mu^+$ then then there exists $G\equiv G'$ with $|G'|=\mu^+$ that contains a clique of cardinality $\mu^+$. 
\end{proposition}
\begin{proof}
By compactness and L\"owenheim-Skolem, it is sufficient to show that $G$ contains arbitrary large finite cliques. Furthermore, by passing to an elementary extension, we may assume that $|G|=2^\mu$. Additionally, after renaming elements, we may assume that $V=2^\mu$. Assume that $\chi(G)\geq \mu^+$.

As $T$ is simple, for every nonzero $\alpha\in V$,  $\tp(\alpha/\{\beta:\beta<\alpha\})$ does not fork over some nonempty countable subset $A_\alpha\subseteq \{\beta:\beta<\alpha\}$. Enumerate $A_\alpha$ as $\langle c_{\alpha,n}:n<\omega\rangle$, possibly with repetitions. Let $F_n$ be the function mapping a nonzero $\alpha\in V$ to $c_{\alpha,n}$ (and $F_n(0)=0$). 

Let $L' \supseteq L\cup \{F_n:n<\omega\}\cup \{<\}$ be a language containing Skolem functions and $G'$ an expansion of $G$ to $L'$ with Skolem functions such that the $F^n$'s are interpreted as above. Let $T'=\mathrm{Th}(G')$. As usual $\Sk(-)$ denotes the Skolem hull in $L'$, i.e. $\Sk(C)$ is the structure generated by $C$ in $G'$. Unless specified otherwise, whatever is done below is done in $L$.  

By forking base monotonicity and the choice of the functions $F_n$, for any $\alpha\in V$, $\tp(\alpha/\{\beta:\beta<\alpha\})$ does not fork over $\Sk(\{\alpha\})\cap \{\beta:\beta<\alpha\}$.
%
%It is clear that:
%\begin{list}{•}{}
%\item for any $a\in M$, $\tp_{\Delta_n}(a/\Sk(\{b:b<a\}))$ is definable over $\Sk(F_n(a))$ and does not fork over $\Sk(F_n(a))$.
%\item for every  $a\in M$, $\tp(a/\Sk(\{b:b<a\}))$ is definable over $\Sk(\bigcup_n F_n(a))$ and does not fork over $\Sk(\bigcup_n F_n(a))$.
%\end{list}
%Note that $\bigcup_n \Sk(F_n(a))=\Sk(\bigcup_n F_n(a))$.

Let $\Delta$ be the collection of all of formulas in one variable $x$ over $\emptyset$ in $L'$ and let $\Delta=\bigcup_{l<\omega}\Delta_l$ be an increasing union of finite subsets. Let $\langle t_i(x): i<\omega\rangle$ be some enumeration of all the terms in $L'$, with $t_0(x)=x$. Thus $\Sk(\{ \alpha\})=\{t_i(\alpha):i<\omega\}$.

Enumerate the $2^\mu$ functions from $\mu$ to $\{0,1\}$ by  $\langle \eta_\alpha\rangle_{\alpha<2^\mu}$, without repetitions.

For any finite subset $u\subseteq \mu$ and $n<\omega$ we define a relation $R_{u,n}$ on 
%\[\dom(R_{u,n}):=\{\alpha\in V=2^\mu: \eta_{t_i(\alpha)}\restriction u\neq \eta_{t_j(\alpha)}\restriction u\text{ for all $i<j<n$ such that $t_i(\alpha)\neq t_j(\alpha)$}\}.\] 
%\begin{gather*}
%\dom(R_{u,n}):=\{\alpha\in V=2^\mu: \eta_{t_i(\alpha)}\restriction u\neq \eta_{t_j(\alpha)}\restriction u\\ \text{ for all $i<j<n$ such that $t_i(\alpha)\neq t_j(\alpha)$}\}. 
%\end{gather*}

\begin{equation*}
\begin{aligned}
\dom(R_{u,n}):=\{\alpha\in V=2^\mu: \eta_{t_i(\alpha)}\restriction u\neq \eta_{t_j(\alpha)}\restriction u \text{ for}\\ \text{all $i<j<n$ such that $t_i(\alpha)\neq t_j(\alpha)$}\}. 
\end{aligned}
\end{equation*}

Let $\alpha \R_{u,n} \beta$ (for $\alpha,\beta\in \dom(R_{u,n})$) if:
\begin{enumerate}
\item for all $j<n$, $\eta_{t_j(\alpha)}\restriction u=\eta_{t_j(\beta)}\restriction u$ and
\item $\tp_{\Delta_n}(\alpha)=\tp_{\Delta_n}(\beta)$.
\end{enumerate}
Note that $R_{u,n}$ is an equivalence relation on $\dom(R_{u,n})$.

\begin{claim}\label{C:exists alpha}
There exists $\alpha\in V$ such that for every finite subset $u\subseteq  \mu$ and $n<\omega$, if $\alpha\in \dom(R_{u,n})$ then 
\[\exists \beta\in [\alpha]_{R_{u,n}} (\alpha\E\beta).\]
\end{claim}
\begin{claimproof}
Note that each $R_{u,n}$ has only finitely many classes on $\dom(R_{u,n})$. Assume towards a contradiction that for every $\alpha\in V$ we can find some $u_\alpha\subseteq \mu$ and $n_\alpha<\omega$ that satisfy the negation of the statement. Now map $\alpha$ to $(u_\alpha,n_\alpha,[\alpha]_{R_{u_\alpha,n_\alpha}})$. This is easily a legal coloring of $V$ by $\mu$ colors, contradicting $\chi(G)\geq \mu^+$.
\end{claimproof}

Let $\alpha\in V$ be as supplied by Claim \ref{C:exists alpha}. For any $n<\omega$ let \[u_n=\{\min\{\varepsilon<\mu: \eta_{t_i(\alpha)}(\varepsilon)\neq \eta_{t_j(\alpha)}(\varepsilon)\}:i<j<n \text{ such that $t_i(\alpha)\neq t_j(\alpha)$}\};\] it is a finite subset of $\mu$. Easily, $\alpha\in \dom(R_{u_n,n})$ for all $n<\omega$. For $n<\omega$ let $\beta_n$ be the element given by Claim \ref{C:exists alpha} for $(u_n,n)$ (and $\alpha$).

Let  $\mathcal{U}$ be a nonprincipal ultrafilter on $\omega$ and let $\widetilde G=(G')^\omega/\mathcal{U}$ be the corresponding ultrapower in the language $L'$. We let $\widetilde G=(\widetilde V,E)$. Set $\widetilde \alpha=[\alpha]_{\mathcal{U}}$, $\widetilde \beta=[\beta_n]_{\mathcal{U}}$; so $\widetilde \alpha \E \widetilde \beta$. 
We make some observations. By definition of $\widetilde \alpha$, $\tp_{L'}(\widetilde \alpha)= \tp_{L'}(\alpha)$ and by definition of the relations $R_{u_n,n}$, this type is also equal to $\tp_{L'}(\widetilde \beta)$. 


Set $\mathbf{a}^{\widetilde \alpha}=\langle t_i(\widetilde \alpha):i<\omega\rangle$ (and likewise for $\alpha$), $\mathbf{a}^{\widetilde \beta}=\langle t_i(\widetilde \beta):i<\omega\rangle$.

\begin{claim}\label{C: enumerate me}
$\mathbf{a}^{\widetilde \alpha}\cap \mathbf{a}^{\widetilde \beta}=\langle t_i(\widetilde \alpha):t_i(\widetilde \alpha)=t_i(\widetilde \beta)\rangle$.
\end{claim}
\begin{claimproof}
Note that for all $i<j<\omega$, if $t_i(\widetilde \alpha)=t_j(\widetilde \alpha)$ then the same holds for $\widetilde \beta$.

We claim that $t_i(\widetilde \alpha)=t_j(\widetilde \beta) \implies t_i(\widetilde \alpha)=t_i(\widetilde \beta)$. Otherwise, assume $i<j$. If $t_i(\widetilde \alpha)=t_j(\widetilde \alpha)$ then, by the first paragraph we are done. So assume not.  We can find some $n$ large enough for which $t_i(\beta_n)\neq t_j(\beta_n)$, $t_i(\alpha)=t_j(\beta_n)$ and $i< j<n$.



As $\beta_n\R_{u_n,n} \alpha$, $\eta_{t_i(\beta_n)}\restriction u_n=\eta_{t_i(\alpha)}\restriction u_n$ so by assumption $\eta_{t_i(\beta_n)}\restriction u_n=\eta_{t_j(\beta_n)}\restriction u_n$, contradicting the definition of $\dom(R_{u_n,n})$.
\end{claimproof}


As $(\mathbf{a}^\alpha,<)$ is a well-order (as a substructure of $2^\mu$) so are $(\mathbf{a}^{\widetilde \alpha},<)$ and $(\mathbf{a}^{\widetilde \beta},<)$. With the aim of applying Lemma \ref{L:technical simple lemma} with $M=\widetilde G$, we prove the following.

\begin{claim}\label{C:forking in tilde G}
For any  $\gamma\in \mathbf{a}^{\widetilde \alpha}\cup\mathbf{a}^{\widetilde \beta}$, $\tp_{L}(\gamma/\{\varepsilon\in \widetilde G :\varepsilon<\gamma\})$ does not fork over $\Sk(\gamma)\cap\{\varepsilon\in\widetilde G:\varepsilon<\gamma\}$.
\end{claim}
\begin{claimproof}
Assume that $\gamma\in \mathbf{a}^{\widetilde \alpha}$. The proof will only use the fact that $\tp_{L'}(\widetilde \alpha)=\tp_{L'}(\alpha)$. Hence, the same proof will also work for $\gamma\in \mathbf{a}^{\widetilde \beta}$. Recall that  $\tp_{L'}(\widetilde{\alpha})=\tp_{L'}(\widetilde \beta)$. Let $t(x)$ be a term (in $L'$) for which $\gamma=t(\widetilde \alpha)$, we get that for $\gamma':=t(\alpha)$, $\tp_{L'}(\gamma)=\tp_{L}(\gamma')$.
%
%\footnote{The reason why $\alpha$ is in parentheses while $\widetilde \alpha$ is not is that the former is an ordinal and the latter is any element, similarly for the elements in the rest of the proof.}

If $\tp_{L}(\gamma/\{\varepsilon :\varepsilon<\gamma\})$ forks over $\Sk(\gamma)\cap\{\varepsilon:\varepsilon<\gamma\}$ then by symmetry of forking there is a tuple $c$ of elements from $\{\varepsilon:\varepsilon <\gamma\}$ such that $\tp_L(c/\{\gamma\}\cup (\Sk(\gamma)\cap \{\varepsilon:\varepsilon<\gamma\}) )$ forks over $\Sk(\gamma)\cap \{\varepsilon:\varepsilon<\gamma\}$. Let $\varphi(y,z)$ be a formula over $\Sk(\gamma)\cap \{\varepsilon:\varepsilon<\gamma\}$, satisfied by $(c,\gamma)$, such that $\varphi(y,\gamma)$  forks over $\Sk(\gamma)\cap \{\varepsilon:\varepsilon<\gamma\}$. Since $\tp_{L'}(\gamma)=\tp_{L}(\gamma')$, $\varphi(y,\gamma')$ forks over $\Sk(\{\gamma'\})\cap \{\varepsilon:\varepsilon<\gamma'\}$.

On the other hand, we know that $\exists y<\gamma \varphi(y,\gamma)$ is in $\tp_{L'}(\gamma/\Sk(\gamma)\cap \{\varepsilon:\varepsilon<\gamma\})$ so $\exists y<\gamma' \varphi(y,\gamma')$ is in $\tp_{L'}(\gamma'/\Sk(\{\gamma'\})\cap \{\varepsilon:\varepsilon<\gamma'\})$. Consequently, there exists a tuple $c'$ of elements in $\{\varepsilon\in V:\varepsilon<\gamma'\}$ for which $\varphi(c',\gamma')$ holds, contradicting the fact that $\tp_L(\gamma'/\{\varepsilon:\varepsilon<\gamma'\})$ does not fork over $\Sk(\{\gamma'\})\cap \{\varepsilon:\varepsilon<\gamma'\}$ (by symmetry of non-forking).
\end{claimproof}

Recall that $\tp_{L'}(\widetilde \alpha)=\tp_{L'}(\widetilde \beta)$. Also, $\mathbf{a}^{\widetilde \alpha}$ and $\mathbf{a}^{\widetilde \beta}$ enumerate elementary substructures  of $\widetilde G\restriction L$. Setting $\mathbf{a}^0=\mathbf{a}^{\widetilde \alpha}\cap \mathbf{a}^{\widetilde\beta}=\langle t_i(\widetilde \alpha):t_i(\widetilde \alpha)=t_i(\widetilde \beta)\rangle$ by Claim \ref{C: enumerate me}, we have $\tp_L(\mathbf{a}^{\widetilde \alpha}/\mathbf{a}^0)=\tp_L(\mathbf{a}^{\widetilde \beta}/\mathbf{a}^0)$. Note that $\mathbf{a}^0$ is an elementary substructure as well.
%
%\begin{claim}
%$\tp^L(\mathbf{a}^\beta/\mathbf{a}^\alpha)$ does not fork over $\mathbf{a}^0$.
%\end{claim}
%\begin{claimproof}
%Assume that $\alpha=[\alpha_n]_\mathcal{U}$ and $\beta=[\beta_n]_\mathcal{U}$ for elements $\alpha_n,\beta_n$ from $C$. By Lemma \ref{L:technical simple lemma}, for any $n<\omega$, $\tp^L(\mathbf{a}^{\beta_n}/\mathbf{a}^{\alpha_n})$ does not fork over $\mathbf{a}^{\alpha_n}\cap\mathbf{a}^{\beta_n}$ (with $L'$ being $L^{sk}$).
%
%Assume, towards a contradiction, that $\tp^L(\mathbf{a}^\beta/\mathbf{a}^\alpha)$ forks over $\mathbf{a}^0$. Hence there is an $L$-formula $\varphi(x_1,\dots,x_l,y_1,\dots,y_m)$ such that  $\varphi(t_{j_1}(x),\dots,t_{j_l}(x), t_{i_1}(\alpha),\dots,t_{i_m}(\alpha))$ is satisfied by $\beta$ and $k$-divides over $\mathbf{a}^0$. Since $k$-dividing is type-definable, there is some $n<\omega$ for which $\varphi(t_{j_1}(x),\dots,t_{j_l}(x), t_{i_1}(\alpha_n),\dots,t_{i_k}(\alpha_n))$ is satisfied by $\beta_n$ and divides over $\mathbf{a}^{\alpha_n}\cap \mathbf{a}^{\beta_n}$; contradiction.
%\end{claimproof}





By Claim \ref{C:forking in tilde G} and Lemma \ref{L:technical simple lemma}, $\mathbf{a}^{\widetilde \alpha}\forkindep[\mathbf{a}^0]\mathbf{a}^{\widetilde \beta}$. So by the independence theorem for simple theories, see  \cite[Lemma 7.4.8]{TZ}, we may find an indiscernible sequence starting with $\mathbf{a}^{\widetilde \alpha}$ and $\mathbf{a}^{\widetilde \beta}$ (in some elementary extension). Thus as $\widetilde \alpha\mathrel{E} \widetilde \beta$ we can find an infinite clique in that elementary extension.
\end{proof}


%
%\begin{proposition}\label{P:cannot embed infinite complete graph}
%Let $T$ be a complete simple theory of graphs in the language of graphs $L=\{E\}$.
%
%Let $\mu\geq 2^{\aleph_0}$ be a cardinal and $G=(V,E)\vDash T$ with $|G|\leq 2^\mu$. If $\chi(G)\geq \mu^+$ then then there exists $G\equiv G'$ with $|G'|=\mu^+$ that contains a clique of cardinality $\mu^+$. 
%\end{proposition}
%\begin{proof}
%By compactness and L\"owenheim-Skolem, it is sufficient to show that $G$ contains arbitrary large finite cliques. Furthermore, by passing to an elementary extension, we may assume that $|G|=2^\mu$. Additionally, after renaming elements, we may assume that $V=2^\mu$. Assume that $\chi(G)\geq \mu^+$.
%
%As $T$ is simple, for every nonzero $\alpha\in V$,  $\tp(\alpha/\{\beta:\beta<\alpha\})$ does not fork over some nonempty countable subset $A_\alpha\subseteq \{\beta:\beta<\alpha\}$. Enumerate $A_\alpha$ as $\langle c_{\alpha,n}:n<\omega\rangle$, possibly with repetitions. Let $F_n$ be the function mapping a nonzero $\alpha\in V$ to $c_{\alpha,n}$ (and $F_n(0)=0$). 
%
%Let $T^{sk}\supseteq T$ be a Skolemization in some language $L^{sk}\supseteq L\cup \{F_n:n<\omega\}$ and let $\Sk(-)$ be the Skolem hull in this language. Unless specified otherwise, whatever is done below is done in $L$.
%
%By forking base monotonicity, for any $\alpha\in V$, $\tp(\alpha/\{\beta:\beta<\alpha\})$ does not fork over $\Sk(\{\alpha\})\cap \{\beta:\beta<\alpha\}$.
%%
%%It is clear that:
%%\begin{list}{•}{}
%%\item for any $a\in M$, $\tp_{\Delta_n}(a/\Sk(\{b:b<a\}))$ is definable over $\Sk(F_n(a))$ and does not fork over $\Sk(F_n(a))$.
%%\item for every  $a\in M$, $\tp(a/\Sk(\{b:b<a\}))$ is definable over $\Sk(\bigcup_n F_n(a))$ and does not fork over $\Sk(\bigcup_n F_n(a))$.
%%\end{list}
%%Note that $\bigcup_n \Sk(F_n(a))=\Sk(\bigcup_n F_n(a))$.
%
%
%We define a relation $R_*$ on $V$ by setting that $\alpha \mathrel{R_*} \beta$ if and only if
%\begin{enumerate}
%\item $\otp(\Sk(\{\alpha\}))=\otp(\Sk(\{\beta\}))$ and let $\pi_{\alpha,\beta}$ be the order isomorphism.
%\item $\pi_{\alpha,\beta}(\alpha)=\beta$.
%\item $\pi_{\alpha,\beta}$ is an elementary partial $L^{sk}\cup\{<\}$-map.
%\end{enumerate}
%
%Note that $R_*$ is an equivalence relation and has to most $2^{\aleph_0}$ equivalence classes. By Fact \ref{F:basic-prop-chi}, there must be an $R_*$-class $C$ such that $\chi(C,E\restriction C)\geq \mu^+$. Note that it is enough that $C$ contains arbitrary large finite cliques. 
%
%By the definition of $R_*$ there is a countable infinite ordinal $\delta$ and a sequence of terms $\langle t_i(x):i<\delta\rangle$, such that for any $\alpha\in C$,  $\mathbf{a}^\alpha=\langle t_i(\alpha):i<\delta\rangle$ is an increasing enumeration of $\Sk(\{\alpha\})$ (without repetitions). 
%
%Let $\widetilde L=L^{sk}\cup\{C,<\}$, i.e. we add a predicate for $C$ and the order relation. Let $\mathcal{U}$ be a nonprincipal ultrafilter on $\omega$ and let $\widetilde G=G^\omega/\mathcal{U}$ be the corresponding ultrapower in the language $\widetilde L$. We let $\widetilde G=(\widetilde V,E)$.
%
%%Naturally, $(\widetilde G,<)$ is no longer a well-order. 
%
%For any $\alpha\in C(\widetilde G)$, note that $\mathbf{a}^\alpha=\langle t_i(\alpha):i<\delta\rangle$ is an increasing  enumeration of $\Sk(\{\alpha\})$ (in $L^{sk}$).
%
%The following claim says that the forking independence property is inherited by $\widetilde G$.
%\begin{claim}\label{C:forking in tilde G}
%For any $\alpha\in C(\widetilde G)$ and $\gamma\in \mathbf{a}^\alpha$, $\tp_{L}(\gamma/\{\beta :\beta<\gamma\})$ does not fork over $\Sk(\gamma)\cap\{\beta:\beta<\gamma\}$.
%\end{claim}
%\begin{claimproof}
%Let $\alpha'\in C(G)$. By the definition of $R_*$,  $\Sk(\alpha)\equiv \Sk(\{\alpha'\})$\footnote{The reason why $\alpha'$ is in parenthesis while $\alpha$ is not is that the former is an ordinal and the latter is any element, similarly for the elements in the rest of the proof.} as two tuples ordered by $<^{\widetilde G}$ (in $L^{sk}\cup \{<\}$). Let $t(x)$ be a term (in $L^{sk}$) for which $\gamma=t(\alpha)$, we get that for $\gamma':=t(\alpha')$, $\Sk(\gamma)\equiv \Sk(\{\gamma'\})$ (in $L^{sk}\cup \{<\}$).
%
%If $\tp_{L}(\gamma/\{\beta :\beta<\gamma\})$ forks over $\Sk(\gamma)\cap\{\beta:\beta<\gamma\}$ then by symmetry there is a tuple $c$ of elements from $\{\beta:\beta <\gamma\}$ such that $\tp_L(c/\{\gamma\}\cup (\Sk(\gamma)\cap \{\beta:\beta<\gamma\}) )$ forks over $\Sk(\gamma)\cap \{\beta:\beta<\gamma\}$. Let $\varphi(x,y)$ be a formula over $\Sk(\gamma)\cap \{\beta:\beta<\gamma\}$, satisfied by $(c,\gamma)$, such that $\varphi(x,\gamma)$  forks over $\Sk(\gamma)\cap \{\beta:\beta<\gamma\}$. Since $\Sk(\gamma)\equiv \Sk(\{\gamma'\})$ in $L^{sk}\cup\{<\}$, $\varphi(x,\gamma')$ forks over $\Sk(\{\gamma'\})\cap \{\beta:\beta<\gamma'\}$.
%
%On the other hand we know that $\exists x<\gamma \varphi(x,\gamma)$ is in $\tp_{L\cup\{<\}}(\gamma/\Sk(\gamma)\cap \{\beta:\beta<\gamma\})$ so $\exists x<\gamma' \varphi(x,\gamma')$ is in $\tp_{L\cup\{<\}}(\gamma'/\Sk(\{\gamma'\})\cap \{\beta:\beta<\gamma'\})$. Consequently, there exists a tuple $c'$ of elements in $\{\beta\in V:\beta<\gamma'\}$ for which $\varphi(c',\gamma')$ holds, contradicting the fact that $\tp_L(\gamma'/\{\beta:\beta<\gamma'\})$ does not fork over $\Sk(\{\gamma'\})\cap \{\beta:\beta<\gamma'\}$ (by symmetry of non-forking).
%\end{claimproof}
%We now note the following.
%\begin{claim}\label{C: enumerate me}
%There exist $\alpha,\beta \in C(\widetilde G)$ such that $\alpha \mathrel{E} \beta$ and that $\mathbf{a}^\alpha\cap \mathbf{a}^\beta=\langle t_i(\alpha):t_i(\alpha)=t_i(\beta)\rangle$.
%\end{claim}
%\begin{claimproof}
%Enumerate the $2^\mu$ functions from $\mu$ to $\{0,1\}$ by  $\langle \eta_\alpha\rangle_{\alpha<2^\mu}$, without repetitions. For any $\alpha,\beta<2^\mu$ let 
%\[H(\alpha,\beta)=
%\begin{cases}
%\min\{\varepsilon<\mu: \eta_\alpha(\varepsilon)=0,\, \eta_\beta(\varepsilon)=1\} +1& \\
%0 & \text{ if the minimum not exist.}
%\end{cases}
%\]
%
%Since $\delta$ is a countable infinite ordinal, there is a bijection  $\sigma:\delta\to \omega$; for any $i<\omega$ set $s_i=t_{\sigma^{-1}(i)}$.
%
%
%For any $n<\omega$, we define a relation $R_n$ on $C$ by setting that $\alpha \mathrel{R_n} \beta$ if and only if for all $i,j<n$, $H(s_i(\alpha),s_j(\alpha))=H(s_i(\beta),s_j(\beta))$.
%Note that each $R_n$ is an equivalence relation and has at most $\mu^{(n^2)}=\mu$ equivalence classes. 
%
%
%If there is some $n<\omega$ for which [$\alpha \mathrel{R_n}\beta$ implies $\neg(\alpha \mathrel{E} \beta)$] then $\chi(C,E\restriction C)=\mu<\mu^+$, contradiction. Thus there is a sequence of elements from $C$, $\langle (\alpha_n,\beta_n):n<\omega\rangle$ satisfying that $\alpha_n\mathrel{R_n} \beta_n$ and $\alpha_n \mathrel{E} \beta_n$.
%
%  Let $\alpha=[\alpha_n]_\mathcal{U}$ and $\beta=[\beta_n]_\mathcal{U}$ be the corresponding elements in the ultrapower; note that both are in $C(\widetilde G)$.
%
%We claim that $s_i(\alpha)=s_j(\beta) \implies i=j$. Otherwise, we can find some $n$ large enough and some $i,j<n$ for which $s_i(\alpha_n)=s_j(\beta_n)$ but $i\neq j$.
%
% As $\alpha_n\mathrel{R_n} \beta_n$, $H(s_i(\alpha_n),s_j(\alpha_n))=H(s_i(\beta_n),s_j(\beta_n))$ and $H(s_j(\alpha_n),s_i(\alpha_n))=H(s_j(\beta_n),s_i(\beta_n))$. Let $H(s_i(\alpha_n),s_j(\alpha_n))=\varepsilon_1$ and $H(s_j(\alpha_n),s_i(\alpha_n))=\varepsilon_2$. 
%
%If $\varepsilon_1>0$ then $0=\eta_{s_i(\alpha_n)}(\varepsilon_1-1)=\eta_{s_j(\beta_n)}(\varepsilon_1-1)=1$, contradiction (where the second equality follows since $s_i(\alpha_n)=s_j(\beta_n)$); likewise it cannot be that $\varepsilon_2>0$, so $\varepsilon_1=\varepsilon_2=0$. Thus \[t_{\sigma^{-1}(i)}(\alpha_n)=s_i(\alpha_n)=s_j(\alpha_n)=t_{\sigma^{-1}(j)}(\alpha_n),\] but as the enumeration of the terms in $\mathbf{a}^{\alpha_n}$ is increasing, this implies that $i=j$, contradiction.
%
%Now, if $t_i(\alpha)\in \mathbf{a}^\beta$ then $t_i(\alpha)=t_i(\beta)$. Indeed, $t_i(\alpha)=t_j(\beta)$ for some $j<\omega$ and thus $s_{\sigma(i)}(\alpha)=s_{\sigma(j)}(\beta)$, so by the above $\sigma(i)=\sigma(j)$ and hence $i=j$.
%\end{claimproof}
%
%
%Let $\alpha,\beta\in C(\widetilde G)$ be elements as supplied by Claim \ref{C: enumerate me}. By the definition of the relation $R_*$ and \L o\'s's theorem, $\mathbf{a}^\alpha\equiv \mathbf{a}^\beta$ as tuples ordered by $<^{\widetilde G}$ (in $L^{sk}\cup\{<\}$). Also, $\mathbf{a}^\alpha$ and $\mathbf{a}^\beta$ enumerate elementary substructures  of $\widetilde G\restriction L$. Setting $\mathbf{a}^0=\mathbf{a}^\alpha\cap \mathbf{a}^\beta=\langle t_i(\alpha):t_i(\alpha)=t_i(\beta)\rangle$, we have $\tp_L(\mathbf{a}^\alpha/\mathbf{a}^0)=\tp_L(\mathbf{a}^\beta/\mathbf{a}^0)$. Note that $\mathbf{a}^0$ is an elementary substructure as well.
%%
%%\begin{claim}
%%$\tp^L(\mathbf{a}^\beta/\mathbf{a}^\alpha)$ does not fork over $\mathbf{a}^0$.
%%\end{claim}
%%\begin{claimproof}
%%Assume that $\alpha=[\alpha_n]_\mathcal{U}$ and $\beta=[\beta_n]_\mathcal{U}$ for elements $\alpha_n,\beta_n$ from $C$. By Lemma \ref{L:technical simple lemma}, for any $n<\omega$, $\tp^L(\mathbf{a}^{\beta_n}/\mathbf{a}^{\alpha_n})$ does not fork over $\mathbf{a}^{\alpha_n}\cap\mathbf{a}^{\beta_n}$ (with $L'$ being $L^{sk}$).
%%
%%Assume, towards a contradiction, that $\tp^L(\mathbf{a}^\beta/\mathbf{a}^\alpha)$ forks over $\mathbf{a}^0$. Hence there is an $L$-formula $\varphi(x_1,\dots,x_l,y_1,\dots,y_m)$ such that  $\varphi(t_{j_1}(x),\dots,t_{j_l}(x), t_{i_1}(\alpha),\dots,t_{i_m}(\alpha))$ is satisfied by $\beta$ and $k$-divides over $\mathbf{a}^0$. Since $k$-dividing is type-definable, there is some $n<\omega$ for which $\varphi(t_{j_1}(x),\dots,t_{j_l}(x), t_{i_1}(\alpha_n),\dots,t_{i_k}(\alpha_n))$ is satisfied by $\beta_n$ and divides over $\mathbf{a}^{\alpha_n}\cap \mathbf{a}^{\beta_n}$; contradiction.
%%\end{claimproof}
%
%
%
%
%
%By Claim \ref{C:forking in tilde G} and Lemma \ref{L:technical simple lemma} (in $L^{sk}$), $\mathbf{a}^\alpha\forkindep[\mathbf{a}^0]\mathbf{a}^\beta$. So by  \cite[Lemma 7.4.8]{TZ} we may find an indiscernible sequence starting with $\mathbf{a}^\alpha$ and $\mathbf{a}^\beta$ (in some elementary extension). Thus as $\alpha\mathrel{E} \beta$ we can find an infinite clique in that elementary extension.
%\end{proof}



%\begin{proof}
%By compactness and L\"owenheim-Skolem, it is sufficient to show that $G$ contains arbitrary large finite cliques. Furthermore, by passing to an elementary extension, we may assume that $|G|=2^\mu$. Additionally, after renaming elements, we may assume that $V=2^\mu$. Assume that $\chi(G)\geq \mu^+$.
%
%As $T$ is simple, for every nonzero $\alpha\in V$,  $\tp(\alpha/\{\beta:\beta<\alpha\})$ does not fork over some nonempty countable subset $A_\alpha\subseteq \{\beta:\beta<\alpha\}$. Enumerate $A_\alpha$ as $\langle c_{\alpha,n}:n<\omega\rangle$, possibly with repetitions. Let $F_n$ be the function mapping a nonzero $\alpha\in V$ to $c_{\alpha,n}$ (and $F_n(0)=0$). 
%
%Let $T^{sk}\supseteq T$ be a Skolemization in some language $L^{sk}\supseteq L\cup \{F_n:n<\omega\}$ and let $\Sk(-)$ be the Skolem hull in this language. Unless specified otherwise, whatever is done below is done in $L$.
%
%By forking base monotonicity, for any $\alpha\in V$, $\tp(\alpha/\{\beta:\beta<\alpha\})$ does not fork over $\Sk(\{\alpha\})\cap \{\beta:\beta<\alpha\})$.
%%
%%It is clear that:
%%\begin{list}{•}{}
%%\item for any $a\in M$, $\tp_{\Delta_n}(a/\Sk(\{b:b<a\}))$ is definable over $\Sk(F_n(a))$ and does not fork over $\Sk(F_n(a))$.
%%\item for every  $a\in M$, $\tp(a/\Sk(\{b:b<a\}))$ is definable over $\Sk(\bigcup_n F_n(a))$ and does not fork over $\Sk(\bigcup_n F_n(a))$.
%%\end{list}
%%Note that $\bigcup_n \Sk(F_n(a))=\Sk(\bigcup_n F_n(a))$.
%
%
%We define a relation $R_*$ on $V$ by setting that $\alpha \mathrel{R_*} \beta$ if and only if
%\begin{enumerate}
%\item $\otp(\Sk(\{\alpha\}))=\otp(\Sk(\{\beta\}))$ and let $\pi_{\alpha,\beta}$ be the order isomorphism.
%\item $\pi_{\alpha,\beta}(\alpha)=\beta$.
%\item $\pi_{\alpha,\beta}$ is an elementary partial $L^{sk}$-map.
%\end{enumerate}
%
%Note that $R_*$ is an equivalence relation and has to most $2^{\aleph_0}$ equivalence classes. By Fact \ref{F:basic-prop-chi}, there must be an $R_*$-class $C$ such that $\chi(C,E\restriction C)\geq \mu^+$. Note that it is enough that $C$ contains arbitrary large finite cliques. 
%
%By the definition of $R_*$ there is a countable (possibly finite) ordinal $\delta$ and a sequence of terms $\langle t_i(x):i<\delta\rangle$, such that for any $\alpha\in C$,  $\mathbf{a}^\alpha=\langle t_i(\alpha):i<\delta\rangle$ is an increasing enumeration of $\Sk(\{\alpha\})$ (without repetitions). 
%
%
%
%
%Enumerate the $2^\mu$ functions from $\mu$ to $\{0,1\}$ by  $\langle \eta_\alpha\rangle_{\alpha<2^\mu}$, without repetitions. For any $\alpha,\beta<2^\mu$ let 
%\[H(\alpha,\beta)=
%\begin{cases}
%\min\{\varepsilon<\mu: \eta_\alpha(\varepsilon)=0,\, \eta_\beta(\varepsilon)=1\} +1& \\
%0 & \text{ if does not exist.}
%\end{cases}
%\]
%
%
%For any $n<\omega$, we define a relation $R_n$ on $C$ by setting that $\alpha \mathrel{R_n} \beta$ if and only if for all $i,j<n$, $H(t_i(\alpha),t_j(\alpha))=H(t_i(\beta),t_j(\beta))$.
%Note that each $R_n$ is an equivalence relation and has at most $\mu^{(n^2)}=\mu$ equivalence classes. 
%
%If there is some $n<\omega$ for which [$\alpha \mathrel{R_n}\beta$ implies $\neg(\alpha \mathrel{E} \beta)$] then $\chi(C,E\restriction C)=\mu<\mu^+$, contradiction. Thus there is a sequence $\langle (\alpha_n,\beta_n):n<\omega\rangle$ satisfying that $\alpha_n\mathrel{R_n} \beta_n$ and $\alpha_n \mathrel{E} \beta_n$.
%
%
%Let $\widetilde L=L^{sk}\cup\{C,<\}$, i.e. we add a predicate for $C$ and the order relation. Let $\mathcal{U}$ be a nonprincipal ultrafilter on $\omega$ and let $\widetilde G=G^\omega/\mathcal{U}$ be the corresponding ultrapower in the language $\widetilde L$. We let $\widetilde G=(\widetilde V,E)$.
%
%Naturally, $(\widetilde G,<)$ is no longer a well-order. For any $\alpha\in C(\widetilde G)$, note that $\mathbf{a}^\alpha=\langle t_i(\alpha):i<\omega\rangle$ is a weakly increasing  enumeration of $\Sk(\{\alpha\})$ (in $L^{sk}$).  Let $\alpha=[\alpha_n]_\mathcal{U}$ and $\beta=[\beta_n]_\mathcal{U}$ be the corresponding elements in the ultrapower; note that both are in $C(\widetilde G)$.
%
%We claim that $t_i(\alpha)=t_j(\beta) \implies t_i(\alpha)=t_j(\alpha)$. Otherwise, we can find some $n$ large enough and some $i,j<n$ for which $t_i(\alpha_n)=t_j(\beta_n)$ but $t_i(\alpha_n)\neq t_j(\alpha_n)$ and thus also $t_i(\beta_n)\neq t_j(\beta_n)$.
%
% As $\alpha_n\mathrel{R_n} \beta_n$, $H(t_i(\alpha_n),t_j(\alpha_n))=H(t_i(\beta_n),t_j(\beta_n))$ and $H(t_j(\alpha_n),t_i(\alpha_n))=H(t_j(\beta_n),t_i(\beta_n))$. Let $H(t_i(\alpha_n),t_j(\alpha_n))=\varepsilon_1$ and $H(t_j(\alpha_n),t_i(\alpha_n))=\varepsilon_2$. 
%
%If $\varepsilon_1>0$ then $0=\eta_{t_i(\alpha_n)}(\varepsilon_1-1)=\eta_{t_j(\beta_n)}(\varepsilon_1-1)=1$, contradiction (where the second equality follows since $t_i(\alpha_n)=t_j(\beta_n)$); likewise it cannot be that $\varepsilon_2>0$. So $\varepsilon_1=\varepsilon_2=0$, but this implies that $t_i(\alpha_n)=t_j(\beta_n)$, contradiction.
%
%This implies that if $t_i(\alpha)\in \mathbf{a}^\beta$ then $t_i(\alpha)=t_i(\beta)$. Indeed, $t_i(\alpha)=t_j(\beta)$ for some $j<\omega$ and thus by the above $t_i(\alpha)=t_j(\alpha)$ so also $t_i(\beta)=t_j(\beta)$.
%
%By the definition of the relation $R_*$ and \L o\'s's theorem, $\tp^{L^{sk}}(\mathbf{a}^\alpha)=\tp^{L^{sk}}(\mathbf{a}^\beta)$. Also, $\mathbf{a}^\alpha$ and $\mathbf{a}^\beta$ enumerate elementary substructures  of $\widetilde G\restriction L$ (in $L$). Setting $\mathbf{a}^0=\mathbf{a}^\alpha\cap \mathbf{a}^\beta=\langle t_i(\alpha):t_i(\alpha)=t_i(\beta)\rangle$, we have $\tp^L(\mathbf{a}^\alpha/\mathbf{a}^0)=\tp^L(\mathbf{a}^\beta/\mathbf{a}^0)$.
%
%\begin{claim}
%$\tp^L(\mathbf{a}^\beta/\mathbf{a}^\alpha)$ does not fork over $\mathbf{a}^0$.
%\end{claim}
%\begin{claimproof}
%For any $\gamma\in C(\widetilde G)$, the induced order $<$ on $\mathbf{a}^\gamma$ is a well-order. As the ultraproduct was taken in the language with the order, it follows that the induced order on $\mathbf{a}^\alpha$ and $\mathbf{a}^\beta$ is a well-order as well. Since a finite union of well orders is a well order, the order induced from $\widetilde G$ on $\Omega:=\mathbf{a}^\alpha\cup \mathbf{a}^\beta$ is still a well-order. We further note that for any $\gamma\in \Omega$, $\Sk(\{\gamma\})\subseteq \Omega$ and  setting $\Omega_{<\gamma}$ for $\{\lambda\in \Omega: \lambda< \gamma\}$,  since it is specified by the type $\{\gamma\}\forkindep[\Sk(\{\gamma\})\cap \Omega_{<\gamma}] \Omega_{<\gamma}$.
%
%
%We will prove by induction on ordinals $\gamma\in \Omega$ that $\mathbf{a}^\alpha\cap \Omega_{<\gamma} \forkindep[\mathbf{a}^\alpha\cap \mathbf{a}^\beta\cap \Omega_{<\gamma}]\mathbf{a}^\beta\cap\ \Omega_{<\gamma}$.
%Since the limit ordinal case is easy, we assume it is know for $\gamma$ and prove it for $\gamma+1$.
%
%%If $\gamma\notin \mathbf{a}^\alpha\cup \mathbf{a}^\beta$ then for $i=1,2$ $\bar a^i\cap (\alpha+1)=\bar a^i\cap (\alpha\cup \{\alpha\})=\bar a^i\cap \alpha$ so there is nothing to show. Otherwise there are three options, either $\alpha\in \bar a^1\setminus \bar a^2$, $\alpha\in \bar a^2\setminus \bar a^1$ or $\alpha\in \bar a^1\cap \bar a^2$.
%
%Assume that $\gamma\in \mathbf{a}^\alpha\setminus \mathbf{a}^\beta$.  We want to prove that $(\mathbf{a}^\alpha\cap \Omega_{<\gamma})\cup\{\gamma\}\forkindep[\mathbf{a}^\alpha\cap \mathbf{a}^\beta\cap \Omega_{<\gamma}] \mathbf{a}^\beta\cap \Omega_{<\gamma}$. By the above, $\{\gamma\}\forkindep[\Sk(\{\gamma\})\cap \Omega_{<\gamma}]\Omega_{<\gamma}$ and thus $\{\gamma\}\forkindep[(\mathbf{a}^\alpha\cap \mathbf{a}^\beta\cap \Omega_{<\gamma})\cup \{\Sk(\{\alpha\})\cap \Omega_{<\gamma}\}] \Omega_{<\gamma}$. As $\gamma\in \mathbf{a}^\alpha$, $\Sk(\{\gamma\})\cap \Omega_{<\gamma} \subseteq \mathbf{a}^\alpha\cap \Omega_{<\gamma}$, $\{\gamma\}\forkindep[(\mathbf{a}^\alpha\cap \mathbf{a}^\beta\cap \Omega_{<\gamma})\cup (\mathbf{a}^\alpha\cap \Omega_{<\gamma})] \Omega_{<\gamma}$ and $\{\gamma\}\forkindep[(\mathbf{a}^\alpha\cap \mathbf{a}^\beta\cap \Omega_{<\gamma})\cup (\mathbf{a}^\alpha\cap \Omega_{<\gamma})] \mathbf{a}^\beta\cap \Omega_{<\gamma}$. By the induction hypothesis and transitivity we conclude.
%
%Assume that $\gamma\in \mathbf{a}^\alpha\cap \mathbf{a}^\beta$. We want to prove that $(\mathbf{a}^\alpha\cap \Omega_{<\gamma})\cup\{\gamma\}\forkindep[(\mathbf{a}^\alpha\cap \mathbf{a}^\beta\cap \Omega_{<\gamma})\cup\{\gamma\}] (\mathbf{a}^\beta\cap \Omega_{<\gamma})\cup\{\gamma\}$. As $\Sk(\{\gamma\})\cap \Omega_{<\gamma}\subseteq \mathbf{a}^\alpha\cap \mathbf{a}^\beta\cap \Omega_{<\gamma}$ and $\{\gamma\}\forkindep[Sk(\{\gamma\})\cap\Omega_{<\gamma}]\Omega_{\gamma}$ we get that $\{\gamma\}\forkindep[\mathbf{a}^\alpha\cap \mathbf{a}^\beta\cap \Omega_{\gamma}] \Omega_{<\gamma}$ and thus $\{\gamma\}\forkindep[(\mathbf{a}^\alpha\cap \mathbf{a}^\beta\cap\Omega_{\gamma})\cup (\mathbf{a}^\alpha\cap \Omega_{<\gamma})] (\mathbf{a}^\beta\cap \Omega_{<\gamma})$. Consequently, $(\mathbf{a}^\alpha\cap \Omega_{<\gamma})\cup\{\gamma\}\forkindep[(\mathbf{a}^\alpha\cap \mathbf{a}^\beta\cap\Omega_{<\gamma})\cup (\mathbf{a}^\alpha\cap \Omega_{\gamma})] (\mathbf{a}^\beta\cap \Omega_{<\gamma})$. As above, it follows by the induction hypothesis and transitivity that $(\mathbf{a}^\alpha\cap \Omega_{<\gamma})\cup\{\gamma\}\forkindep[(\mathbf{a}^\alpha\cap \mathbf{a}^\beta\cap\Omega_{<\gamma})] (\mathbf{a}^\beta\cap \Omega_{<\gamma})$. We can now conclude. 
%\end{claimproof}
%
%
%
%
%
%In particular, all the above implies that $\tp^L(\mathbf{a}^\alpha/\mathbf{a}^0)=\tp^L(\mathbf{a}^\beta/\mathbf{a}^0)$ and $\mathbf{a}^\alpha\forkindep[\mathbf{a}^0]\mathbf{a}^\beta$ and consequently by \cite[Lemma 7.4.8]{TZ} we may find an indiscernible sequence starting with $\mathbf{a}^\alpha$ and $\mathbf{a}^\beta$ (in some elementary extension). Thus as $\alpha\mathrel{E} \beta$ we can find an infinite clique in that elementary extension.
%\end{proof}
%
%






Given Proposition \ref{P:cannot embed infinite complete graph}, a natural question to ask if we can necessarily find an infinite clique already in $G$ it self. The following example shows that you cannot hope for an uncountable clique in general.

Recall that every random graph has a simple theory and that every such graph contains a countable infinite clique so these cannot be avoided.

\begin{proposition}\label{P:embed into random}
Let $\mu$ be an infinite cardinal. Any graph of cardinality $\mu$ without an uncountable clique can be embedded in a random graph of cardinality $\mu$ with no uncountable clique.
%
%For any cardinal $\mu$ there exists a random graph of cardinality and chromatic number $\mu$ with no uncountable clique. 
\end{proposition}
\begin{proof}
Let $\mu$ be an infinite cardinal and let $G_0=(V_0,E_0)$ be the given graph. Without loss of generality, $V_0=\{2\alpha: \alpha<\mu\}$.


% be any graph of cardinality $\mu$ with no infinite clique and $\chi(G_0)=\mu$ (for example the triangle-free graph from \cite{ErRa-triangle-free} \textbf{move to introduction?}). Without loss of generality, $V_0=\{2\alpha: \alpha<\mu\}$.

Let $\langle u_\gamma\rangle_{\gamma<\mu}$ enumerate all finite subsets of $\mu$ such that  each finite subset  occurs $\mu$ times. In particular, for any $\gamma<\alpha<\mu$ there is some $\alpha<\gamma'<\mu$ for which $u_\gamma=u_{\gamma'}$.

Let $V=\mu$. We define a new graph $G=(V,E)$, extending $G_0$, such that for $\alpha< \beta<\mu$ if $\beta=2\gamma+1$ for some $\gamma<\mu$ and $\alpha\in u_\gamma$ we let $\{\alpha,\beta\}$ be an edge.

We claim that $G$ has the desired properties. 
%Since $G_0$ embeds into $G$ and $|G|=\mu$ it follows that $\chi(G)=\mu$.

We note that $G$ is random graph. Indeed, let $X=\{\alpha_1,\dots,\alpha_n\},\,Y=\{\beta_1,\dots,\beta_n\}$ be two disjoint sets of vertices. Let $\gamma<\mu$ be  larger than the $\alpha_i$'s and $\beta_j$'s and satisfying that $u_\gamma=\{\alpha_1,\dots, \alpha_n\}$. Then $2\gamma+1$ is connected to each of the $\alpha_i$ and to none of the $\beta_j$.

Finally, assume that $G$ contains a clique $C$ of cardinality $\aleph_1$. By assumption, $C\cap V_0$ must be at most countable so there must a clique of cardinality $\aleph_1$ consisting of odd ordinals in $\mu$. Let $U$ be the first $\omega$ of those and let $2\gamma+1\in C$ be an ordinal larger than any of the ordinals in $U$. But $2\gamma+1$ can only be connected to finitely many vertices which are smaller, contradiction. 
\end{proof}

\begin{corollary}
For any infinite cardinal $\mu$ there exists a random graph of cardinality and chromatic number $\mu$ with no uncountable clique. 
\end{corollary}
\begin{proof}
Apply Proposition \ref{P:embed into random} to any graph $G_0$ of cardinality $\mu$ with no infinite clique and $\chi(G_0)=\mu$ (for example the triangle-free graph from \cite{ErRa-triangle-free}).

Let $G$ be the graph supplied by the proposition. Since $G_0$ embeds into $G$ and $|G|=\mu$ it follows that $\chi(G)=\mu$.
\end{proof}


\begin{question}
Is there a theory of simple graphs such that for every cardinal $\mu$ we can find a graph of cardinality and chromatic number $\mu$ with no infinite cliques at all?
\end{question}


\begin{question}
Does an analog of Proposition \ref{P:cannot embed infinite complete graph} hold  for other model theoretic tame graphs, such as NSOP$_1$ and NIP?
\end{question}

%
%\begin{lemma}\label{L:cannot embed infinite complete graph}
%Let $\mu=\mu^{\aleph_0}$ and $G\vDash T$ with $|G|\leq 2^\mu$. If $G$ does not contain arbitrary finite cliques then $\chi(G)\leq \mu$.
%\end{lemma}
%\begin{proof}
%By replacing $G$ with an elementary extension, we may assume that $|G|=2^\mu$. Furthermore, after renaming elements, we assume that $G=2^\mu$.
%
%For every $a\in G$, every $p\in S(\{b:b<a\})$ does not fork over some countable subset of $\{b:b<a\}$. Assume that $\tp(a/\{b:b<a\})$ does not fork over $\{c_{a,n}:n<\omega\}\subseteq \{b:b<a\}$. For any $a\in G$ let $F_n$ be the function mapping $a$ to $c_{a,n}$. 
%
%Let $T\subseteq T^{sk}$ be a Skolemization in some language $L\cup \{F_n\}_{n<\omega}\subseteq L^{sk}$ and let $\Sk(-)$ be the Skolem hull in this language. Unless specified otherwise, whatever is done below is done in $L$.
%
%By properties of forking, for any $a\in G$, $\tp(a/\{b:b<a\})$ does not fork over $\Sk(\{a\})\cap \{b:b<a\})$.
%%
%%It is clear that:
%%\begin{list}{•}{}
%%\item for any $a\in M$, $\tp_{\Delta_n}(a/\Sk(\{b:b<a\}))$ is definable over $\Sk(F_n(a))$ and does not fork over $\Sk(F_n(a))$.
%%\item for every  $a\in M$, $\tp(a/\Sk(\{b:b<a\}))$ is definable over $\Sk(\bigcup_n F_n(a))$ and does not fork over $\Sk(\bigcup_n F_n(a))$.
%%\end{list}
%%Note that $\bigcup_n \Sk(F_n(a))=\Sk(\bigcup_n F_n(a))$.
%
%Choose some distinct $\langle \eta_\alpha:\mu \to \{0,1\} : \alpha<2^\mu\rangle$. 
%For any $\alpha,\beta<2^\mu$ let 
%\[H(\alpha,\beta)=
%\begin{cases}
%\min\{\varepsilon<\mu: \eta_\alpha(\varepsilon)=0,\, \eta_\beta(\varepsilon)=1\} +1& \\
%0 & \text{ if does not exist.}
%\end{cases}
%\]
%
%We define a relation $R$ on $2^\mu$ by setting that $\alpha \mathrel{R} \beta$ if and only if
%\begin{enumerate}
%\item $\otp(\Sk(\{\alpha\}))=\otp(\Sk(\{\beta\}))$ and let $\pi_{\alpha,\beta}$ be the order isomorphism
%\item $\pi_{\alpha,\beta}(\alpha)=\beta$.
%\item $\pi_{\alpha,\beta}$ is a partial $L$-isomorphism 
%\item For any $\zeta<\mu$, $\pi_{\alpha,\beta}$ preserves the relation $H(-,-)=\zeta$.
%\end{enumerate}
%Note that $R$ is an equivalence relation and has at most $\mu^{\aleph_0}=\mu$ equivalence classes. Our aim is to prove that if $\alpha\R \beta$ then $\neg \alpha\E\beta$. That will certainly be enough for then we can color the vertices of the graph by the using the different $R$-classes.
%
%Let $\alpha_1,\alpha_2<2^\mu$ be such that $\alpha_1\R \alpha_2$.
%
%For $l=1,2$, let $\overline a^l=\langle a^l_i: i<\lambda\rangle$ be an increasing enumeration of $\Sk(\{\alpha_l\})$. By the definition of the relation $R$, $\tp^L(\overline a^1)=\tp^L(\overline a^2)$ and by the existence of Skolem functions, $\overline a^1$ and $\overline a^2$ are elementary submodels (in $L$).
%
%
%We claim that $\overline a^0:=\overline a^1\cap \overline a^2=\langle a^1_i: a^1_i=a^2_i\rangle$. Indeed, assume that $a^1_i=a^2_j$ with $i\neq j$. As $\alpha_1\mathrel{R} \alpha_2$, $H(a^1_i,a^1_j)=H(a^2_i,a^2_j)$ and $H(a^1_j,a^1_i)=H(a^2_j,a^2_i)$. Let $H(a^1_i,a^1_j)=\varepsilon_1$ and $H(a^1_j,a^1_i)=\varepsilon_2$. 
%
%If $\varepsilon_1>0$ then $0=\eta_{a_i^1}(\varepsilon_1-1)=\eta_{a^2_j}(\varepsilon_1-1)=1$, contradiction; likewise it cannot be that $\varepsilon_2>0$. So $\varepsilon_1=\varepsilon_2=0$, but this implies that $\eta$, contradiction.
%
%By the existence of Skolem functions, $\bar a^0$ is also an elementary substructure. Since $\tp^L(\bar a^1)=\tp^L(\bar a^2)$ it follows that by the above that $\tp^L(\bar a^1/\bar a^0)=\tp^L(\bar a^2/\bar a^0)$.
%
%\begin{claim}
%$\tp^L(\overline a^2/\overline a^1)$ does not fork over $\overline a^0$.
%\end{claim}
%\begin{proof}
%We will prove by induction on ordinals $\alpha\in G$ that $\bar a^1\cap \alpha \forkindep[\bar a^1\cap \bar a^2\cap \alpha]\bar a^2\cap \alpha$.
%Since the limit ordinal case is easy, we assume it is know for $\alpha$ and prove it for $\alpha+1$.
%
%If $\alpha\notin \bar a^1\cup \bar a^2$ then for $i=1,2$ $\bar a^i\cap (\alpha+1)=\bar a^i\cap (\alpha\cup \{\alpha\})=\bar a^i\cap \alpha$ so there is nothing to show. Otherwise there are three options, either $\alpha\in \bar a^1\setminus \bar a^2$, $\alpha\in \bar a^2\setminus \bar a^1$ or $\alpha\in \bar a^1\cap \bar a^2$.
%
%Assume that $\alpha\in \bar a^1\setminus \bar a^2$.  We want to prove that $(\bar a^1\cap \alpha)\cup\{\alpha\}\forkindep[\bar a^1\cap\bar a^2\cap \alpha] \bar a^2\cap \alpha$. By the above, $\{\alpha\}\forkindep[\Sk(\{\alpha\})\cap \alpha]\alpha$ and thus $\{\alpha\}\forkindep[(\bar a^1\cap\bar a^2\cap \alpha)\cup \{\Sk(\{\alpha\})\cap \alpha\}] \alpha$. As $\alpha\in \bar a^1$, $\Sk(\{\alpha\})\cap \alpha \subseteq \bar a^1\cap \alpha$, $\{\alpha\}\forkindep[(\bar a^1\cap\bar a^2\cap \alpha)\cup (\bar a^1\cap \alpha)] \alpha$ and $\{\alpha\}\forkindep[(\bar a^1\cap\bar a^2\cap \alpha)\cup (\bar a^1\cap \alpha)] \bar a^2\cap \alpha$. By the induction hypothesis and transitivity we conclude.
%
%Assume that $\alpha\in \bar a^1\cap \bar a^2$. We want to prove that $(\bar a^1\cap \alpha)\cup\{\alpha\}\forkindep[(\bar a^1\cap\bar a^2\cap \alpha)\cup\{\alpha\}] (\bar a^2\cap \alpha)\cup\{\alpha\}$. As $\Sk(\{\alpha\})\cap \alpha\subseteq \bar a^1\cap \bar a^2\cap \alpha$ and $\{\alpha\}\forkindep[Sk(\{\alpha\})\cap\alpha]\alpha$ we get that $\{\alpha\}\forkindep[\bar a^1\cap \bar a^2\cap\alpha] \alpha$ and thus $\{\alpha\}\forkindep[(\bar a^1\cap \bar a^2\cap\alpha)\cup (\bar a^1\cap \alpha)] (\bar a^2\cap \alpha)$. Consequently, $(\bar a^1\cap \alpha)\cup\{\alpha\}\forkindep[(\bar a^1\cap \bar a^2\cap\alpha)\cup (\bar a^1\cap \alpha)] (\bar a^2\cap \alpha)$. As above, it follows by the induction hypothesis and transitivity that $(\bar a^1\cap \alpha)\cup\{\alpha\}\forkindep[(\bar a^1\cap \bar a^2\cap\alpha)] (\bar a^2\cap \alpha)$. We can now conclude. 
%\end{proof}


\subsection{Graphs with Stable Edge Relation}
Before getting into the main result we prove a technical lemma which may be interesting on its own.


\begin{lemma}\label{L:club of models}
Let $T$ be a first order theory, $\mu$ an infinite cardinal and $\varphi(x,y)$ a stable formula. Let $M\vDash T$ with $|M|=\mu^+$ and assume that $M$ is an increasing continuous union of elementary substructures $\langle M_\alpha\rangle_{\alpha<\mu^+}$ each of cardinality at most $\mu$.

Let $\psi(y,z)$ be a uniform definition of $\varphi$-types (as in Fact \ref{F:def of types}). For any $a\in M$ and $\alpha<\mu^+$, let $c_{a,\alpha}\in M_\alpha$ be such that $\psi(y,c_{a,\alpha})$ defines $\tp_\varphi(a/M_\alpha)$.

Then there exists a club $\mathcal{C}\subseteq \mu$ of limit ordinals satisfying that for any $\delta\in \mathcal{C}$ and $a\in M\setminus M_\delta$:

$(\dagger)_{a,\delta}$ For any $\delta<\beta<\mu^+$ there is $b\in M\setminus M_\beta$ for which $\tp_\varphi(b/M_\beta)$ is definable by $\psi(y,c_{a,\delta})$.
\end{lemma}
\begin{proof}
We first prove that  for any $\alpha<\mu^+$, $(\dagger')_\alpha$ there is some limit ordinal $\delta>\alpha$ such $(\dagger)_{a,\delta}$ holds for any $a\in M\setminus M_\delta$. Assume otherwise, i.e. there exists an $\alpha<\mu^+$ for which for any limit ordinal $\delta>\alpha$ there is some $a_\delta\in M\setminus M_\delta$ and $\beta_\delta>\delta$ such that for any $b\in M\setminus M_{\beta_\delta}$, $\tp_\varphi(b/M_{\beta_\delta})$ is not defined by $\psi(y,c_{a_\delta,\delta})$.

For any limit ordinal $\delta>\alpha$, let $f(\delta)$ be the minimal ordinal $\varepsilon$ for which $c_{a_\delta,\delta}\in M_\varepsilon$. Note that as $c_{a_\delta,\delta}$ is a finite tuple and $\delta$ is a limit ordinal, necessarily $f(\delta)<\delta$. As the set of limit ordinals between $\alpha$ and $\mu^+$ is a stationary subset of $\mu^+$ (it is even a club), by Fodor's lemma (\cite[Theorem 8.7]{jech}) there exists a stationary subset $S\subseteq \mu^+$ and $\varepsilon<\mu^+$ for which $f(\delta)=\varepsilon$ for any $\delta\in S$.

By definition, $c_{a_\delta,\delta}\in M_\varepsilon$ for any $ \delta\in S$. As $|M_\varepsilon|\leq \mu$, by the pigeonhole principle there is an unbounded subset $S'\subseteq S$ (i.e. of cardinality $\mu^+$) for which $c:=c_{a_{\delta_1},\delta_1}=c_{a_{\delta_2},\delta_2}$ for any $\delta_1,\delta_2\in S'$.

Now, pick any $\delta\in S'$. By our assumption there is some $\beta_\delta>\delta$ such that for any $b\in M\setminus M_{\beta_\delta}$, $\tp_\varphi(b/M_{\beta_\delta})$ is not defined by $\psi(y,c_{a_\delta,\delta})=\psi(y,c)$.

Let $\beta'>\beta_\delta$ be an element in $S'$ (it is unbounded) and let $a_{\beta'}$ be the corresponding element. So $a_{\beta'}\notin M_{\beta'}$ and in particular $\notin  M_{\beta_\delta}$ and thus $\tp_\varphi(a_{\beta'}/M_{\beta_\delta})$ is not defined by $\psi(y,c_{a_\delta,\delta})=\psi(y,c)$. On the other hand, by choice of $S'$, $\tp_\varphi(a_{\beta'}/M_{\beta'})$ and thus also $\tp_\varphi(a_{\beta'}/M_\beta)$ is defined by $\psi(y,c)$, contradiction. 

Now, we turn to show that we can find such a club $\mathcal{C}$. Let $g:\mu^+\to \mu^+$ be the function mapping $\alpha$ to the minimal limit ordinal $\delta$ satisfying $(\dagger')_\alpha$ and let $\mathcal{C}=\{\alpha<\delta<\mu^+: \text{ $\delta$ is a limit ordinal and } \forall(\alpha<\delta) g(\alpha)<\delta\}$; it is a club in $\mu^+$.

Let $\delta\in \mathcal{C}$ and $a\in M\setminus M_\delta$. As $\delta$ is a limit ordinal, $c_{a,\delta}\subseteq M_\varepsilon$ for some $\varepsilon<\delta$; so $g(\varepsilon)<\delta$. Let $\delta<\beta<\mu^+$; we need to show that there exists $b\in M\setminus M_\beta$ for which $\tp_\varphi(b/M_\beta)$ is defined by $\psi(y,c_{a,\delta})$.

By $(\dagger)_{a,g(\varepsilon)}$, since $a\notin M_{g(\varepsilon)}$ and $\beta>\delta>g(\varepsilon)$, there is $b\in M\setminus M_\beta$ for which  $\tp_\varphi(b/M_\beta)$ is defined by $\psi(y,c_{a,\delta})$, as required.
% we have that for any $\beta'>\beta$ there is some $b_{\beta'}\notin M_{\beta'}$ for which $\tp_\varphi(b_{\beta'}/M_{\beta'})$ is not defined by by $\psi(y,c_{a,\delta})$.
\end{proof}


We phrase forking symmetry for a stable formula in a form useful to us.
\begin{fact}[Harrington]\label{F:local Harrington}
Let $T$ be a first order theory with monster model $\mathbb{U}$. Let $\varphi(x,y)$ be a stable formula and $\psi(y,z)$ a formula uniformly defining $\varphi$-types. For any two small models $N_1$ and $N_2$ and elements $a,b\in \mathbb{U}$, if $\psi(y,c_a)$ defines $\tp_\varphi (a/N_1)$ and $\psi(y,c_b)$ defines $\tp_\varphi(b/N_2)$ and $c_a\in N_2,\, c_b\in N_1$ then 
\[\psi(a,c_b)\Leftrightarrow \psi(b,c_a).\]
\end{fact}
\begin{proof}
Let $p=\tp_\varphi(a/N_1)$ and $q=\tp_\varphi(b/N_2)$.
Let $\widetilde p\supseteq p$ be the global $\varphi$-type that $\psi(y,c_a)$ defines and $\widetilde q\supseteq q$ be the global $\varphi$-type that $\psi(y,c_b)$ defines. By \cite[Lemma 8.3.4]{TZ}, $\psi(x,c_b)\in \widetilde p\Leftrightarrow \psi(x,c_a)\in \widetilde q$; but as $c_a\in N_2$ and $c_b\in N_1$, we conclude.
\end{proof}

The following is due to Engelking-Kar\l owicz \cite{En-Ka}, see also \cite{Rinot}.
\begin{fact}\label{F:En-Ka}
For cardinals $\kappa\leq \lambda\leq \mu\leq 2^\lambda$ the following are equivalent:
\begin{enumerate}
\item $\lambda^{<\kappa}=\lambda$
\item there exists a collection of functions $\langle f_i:\mu\to \lambda\rangle_{i<\lambda}$, such that for every $X\in [\mu]^{<\kappa}$ and every function $f:X\to \lambda$, there exists some $i<\lambda$ with $f\subseteq f_i$.
\end{enumerate}
\end{fact}

We now prove the main result of this section.

\begin{proposition}\label{P:stable}
Let $L=\{E\}$ be the language of graphs and $T$ be an $L$-theory specifying that $E$ is a symmetric and irreflexive stable relation. For any infinite cardinal $\mu$ and $G\vDash T$ with $|G|=\mu^+$, if $\chi(G)\geq \mu^+$ then $G$ contains an infinite clique of cardinality $\mu^+$.
\end{proposition}
\begin{proof}
For ease of notation, write $\varphi(x,y)=E(x,y)$ and let $\psi(y,z)$ be a uniform definition for $\varphi$-types (as in Fact \ref{F:def of types}). Assume that $\chi(G)\geq \mu^+$.

By Lemma \ref{L:club of models}, there exists an increasing  continuous family of elementary substructures $\{G_\alpha\prec G:0<\alpha<\mu^+\}$   with $|G_\alpha|=\mu$ and $\bigcup_{0<\alpha<\mu^+} G_\alpha=G$, satisfying:

$(\dagger)$ For any $0<\delta<\mu^+$ and $a\in G\setminus G_\delta$  there exists $c_{a,\delta}\in G_\delta$ such that $\psi(y,c_{a,\delta})$ defines $\tp_\varphi(a/G_\delta)$ satisfying that  for any $\delta<\beta<\mu^+$ there is $b\notin G_\beta$ for which $\tp_\varphi(b/G_\beta)$ is defined by $\psi(y,c_{a,\delta})$. 

Set $G_0=\emptyset$

For any $a\in G$ let $\alpha^a_0<\mu^+$ be minimal such that $a\in G_{\alpha_0+1}$. Let $c_0^a=c_{a,\alpha_0^a}\in G_{\alpha^a_0}$, in particular $\psi(y,c_0^a)$ defines $\tp_\varphi (a/G_{\alpha^a_0})$. Let $\alpha_1^a<\alpha_0^a$ be such that $c_0^a\in G_{\alpha_1^a+1}\setminus G_{\alpha^a_1}$. Let $c^a_1=c_{a,\alpha_1^a}$. Likewise we continue and find a sequence $\langle (c_{k-1}^a,\alpha_k^a):1\leq k\leq n_a\rangle$ satisfying
\begin{list}{•}{}
\item $c^a_{k-1}=c_{a,\alpha^a_{k-1}}$, in particular $\psi(y,c^a_{k-1})$ defines $\tp_\varphi(a/G_{\alpha^a_{k-1}})$
\item $c^a_{k-1}\in G_{\alpha^a_k+1}\setminus G_{\alpha^a_k}$. 
\item $\alpha^a_{n_a}=0$.
\end{list}

Note that $\alpha_0=0$ if and only if $n_a=0$ and then there are no $c$'s.

We now define a coloring.

By  Engelking-Kar\l owicz (Fact \ref{F:En-Ka}), there exists a family of functions $\{g_{\beta}:\mu^+\to \mu:\beta<\mu\}$ satisfying that for any finite subset $X\subseteq \mu^+$ and every function $f:X\to \mu$ there is some function $g_\beta$, $\beta<\mu$, with $f\subseteq g_\beta$.


For any such $a\in G$ define a function $f_a:\{\alpha^a_0,\dots,\alpha_{n_a}^a\}\to \omega$ by setting $f_a(\alpha^a_i)=i$. Let $\beta^a<\mu$ be minimal such that $f_a\subseteq g_{\beta^a}$.

For any $\alpha<\mu^+$ let $s_\alpha :G_\alpha\to \mu$ be a bijection. For simplicity, we also denote by $s_\alpha$ the induced bijection between $G_\alpha^n$ and $\mu^n$. 
Let $Y=\aleph_0 \times \mu^2\times \mu^{<\omega}$; note that $|Y|=\mu$.
%
%Further, for any $n<\omega$ we choose a bijection $\Psi_n$ between $\mu^{n}$ and $\mu$.

Let $\phi:G\to Y$ be a function mapping $a\in G$ to \[(n_a, \beta^a,s_{\alpha^a_0+1}(a),(s_{\alpha_i^a+1}(c^a_{i-1}))_{1\leq i\leq n_a}).\] Since $\phi$ cannot be a legal coloring there exist $a,b\in G$ distinct elements satisfying that $\phi(a)=\phi(b)$ and $\varphi(a,b)$ holds, i.e. $a\mathrel{E} b$. 


%
%For any $\alpha<\mu^+$ let $s_\alpha :G_\alpha\to \mu$ be a bijection. For simplicity, we also denote by $s_\alpha$ the induced bijection between $G_\alpha^n$ and $\mu^n$. 
%Let $S^{\text{fin}}_\varphi(G_\alpha)$ be the collection of all $\varphi$-types over finite subsets of $G_\alpha$. Note that $|S^{\text{fin}}_\varphi(G_\alpha)|=\mu$.
%Let $s_\alpha':S^{\text{fin}}_\varphi(G_\alpha)\to \mu$ be a bijection. Let $Y=\aleph_0 \times \mu^2\times \mu^{<\omega}\times \mu^{<\omega}$; note that $|Y|=\mu$.
%%
%%Further, for any $n<\omega$ we choose a bijection $\Psi_n$ between $\mu^{n}$ and $\mu$.
%
%Let $\phi:G\to Y$ be a function mapping $a\in G$ to \[(n_a, \beta^a,s_{\alpha^a_0+1}(a),(s_{\alpha_i^a+1}(c^a_{i-1}))_{1\leq i\leq n_a},(s_\alpha'(\tp_\varphi(a/c^a_{i-1})))_{1\leq i\leq n_a}).\] Since $\phi$ cannot be a legal coloring there exist $a,b\in G$ distinct elements satisfying that $\phi(a)=\phi(b)$ and $\varphi(a,b)$ holds, i.e. $a\mathrel{E} b$. 


Without loss of generality, assume that $\alpha^a_0\geq \alpha^b_0$. Note that necessarily, $\alpha^a_0>\alpha^b_0$, since otherwise, as $s_{\alpha_0^a+1}(a)=s_{\alpha_0^b+1}(b)$ we conclude that $a=b$, contradiction. In particular, $n_a>0$ and so also $n_b=n_a>0$.

Also, if $\alpha_i^a=\alpha^b_j$ for some $i,j$ then, since by assumption $\beta^a=\beta^b$, $i=j$.


Let $n>0$ be minimal such that $\alpha_n^a=\alpha^b_n$ (such $n$ exists since this equality holds for $n=n_a>0$). Consider the (ordered) set $A'=\{\alpha_i^a,\alpha_i^b:i< n\}$. Note that the elements in $A'$ are distinct. Now let $A=A'\cup\{\alpha^a_n=\alpha^b_n\}$.

We will call an element $\alpha_i^a\in A$ (with $i>0$) an \emph{$a$-pivot} if there exists an element $\alpha_j^b\in A$ with $\alpha^a_{i-1}>\alpha_j^b>\alpha_i^a$ (and likewise a \emph{$b$-pivot} for $\alpha_i^b\in A$ with $i>0$). Note that $\alpha^a_n=\alpha^b_n$ is either an a-pivot or a b-pivot.

We will prove the following by (downward) induction.
\begin{claim}
For every $a$-pivot $\alpha_i^a\in A$, with $i>0$, $\psi(b,c^a_{i-1})$ holds and for every $b$-pivot $\alpha_i^b\in A$, with $i>0$, $\psi(a,c^b_{i-1})$ holds.
\end{claim}
\begin{claimproof}
%We first two preliminary  results:
%
%Let $l_1<n$ be maximal such that $b\in M_{\alpha^a_{l_1}}$ (and hence $\alpha^b_0<\alpha^a_{l_1}$). As $\varphi(a,b)$, we have that $\varphi(x,b)\in \tp_\varphi(a/M_{\alpha^a_{l_1}})$ and so $\psi(b,c_{\alpha^a_{l_1}})$.
%
%Let $r_1<n$ be maximal with $\alpha_{r_1}^b>\alpha^a_{l_1+1}$. We have that $c^a_{l_1}\in M_{\alpha^a_{l_1+1}+1}\subseteq M_{\alpha^b_{r_1}}$ and so by the above $\psi(y,c^a_{l_1})\in \tp_\varphi(b/M_{\alpha^b_{r_1}})$. Since $c^b_{r_1-1}\in M_{\alpha^b_{r_1}+1}\subseteq M_{\alpha^a_{l_1}}$,  by symmetry (cite), we get that $\psi(y,c^b_{r_1-1})\in \tp_\varphi (a/M_{\alpha^a_{l_1}})$ as well.
%
%Now we treat the pivots. 
Let $\alpha^a_{t+1}\in A$ be an $a$-pivot (possibly $t+1=n$). Let $v<n$ be minimal with $\alpha^a_t>\alpha^b_{v}>\alpha^a_{t+1}$; $\alpha^b_v$ is either a $b$-pivot and $v=s+1$ for some $s$ or $v=0$. Assume, first, that $\alpha^b_{s+1}$ is a $b$-pivot; so $\psi(a,c^b_s)$ holds by the induction hypothesis. As $c^b_s\in G_{\alpha^b_{s+1}+1}\subseteq G_{\alpha_t^a}$ and $c^a_t\in G_{\alpha^a_{t+1}+1}\subseteq G_{\alpha^b_{s+1}}\subseteq G_{\alpha^b_s}$, we can apply Fact \ref{F:local Harrington} with $\tp_\varphi(a/G_{\alpha^a_t})$ and $\tp_\varphi(b/G_{\alpha^b_s})$ and conclude that $\psi(b,c^a_t)$ holds.  

Now assume that $v=0$. Since $b\in G_{\alpha^b_0+1}\subseteq G_{\alpha^a_t}$ and $\varphi(a,b)$ holds then, as $\psi(y,c^a_t)$ defines $\tp_\varphi (a/G_{\alpha^a_t})$ we get that $\psi(b,c^a_t)$ holds; as needed.

The case where $\alpha^b_{t+1}\in A$ is a $b$-pivot is proved similarly.
\end{claimproof}

Note that $c^a_{n-1}\in G_{\alpha^a_n+1}=G_{\alpha^b_n+1}\ni c^b_{n-1}$. As $\phi(a)=\phi(b)$, we necessarily have $c:=c^a_{n-1}=c^b_{n-1}$.
% and $a\equiv_c^\varphi b$. 
 If $\alpha^a_n$ is an a-pivot we have that $\psi(a,c)$ holds and if it is a b-pivot then $\psi(b,c)$ holds. Assume the former holds (the proof where the latter holds is identical).
 
% Either way, since $\psi$ is a boolean combination of $\varphi^{\text{opp}}$ formulas, we have that both $\psi(a,c)$ and $\psi(b,c)$ hold.
%
%
%Recall that $a\in G_{\alpha_0^a+1}\setminus G_{\alpha_0^a}$,  $b\in G_{\alpha^b_0+1}\setminus G_{\alpha_0^b}$. 

We inductively construct a sequence $(d_\alpha)_{\alpha<\mu^+}$ in $V$ such that:
\begin{list}{•}{}
\item $\psi(d_\alpha,c)$ holds for all $\alpha<\mu^+$.
\item $\varphi(d_\alpha,d_\beta)$ for all $\alpha\neq \beta<\mu^+$.
\end{list}

Suppose we constructed $(d_\alpha)_{\alpha<\gamma}$. Let $\alpha_{n-1}^a<\delta<\mu^+$ be such that $d_\alpha\in G_\delta$ for any $\alpha<\gamma$. By $(\dagger)$ there exists $d_\gamma\in G\setminus G_\delta$ such that $\tp_\varphi(d_\gamma/G_\delta)$ is definable by $\psi(x,c)$.

Since $\psi(d_\alpha,c)$ holds for any $\alpha<\gamma$ it follows that $\varphi(d_\gamma,d_\alpha)$ holds (and also $\varphi(d_\alpha,d_\gamma)$ by symmetry). Additionally, $\psi(x,c)$ defines both $\tp_\varphi(a/G_{\alpha^a_{n-1}})$ and $\tp_\varphi(d_\gamma/G_{\alpha^a_{n-1}})$; hence $a\equiv^\varphi_{G_{\alpha^a_{n-1}}} d_\gamma$.  As $c\in G_{\alpha^a_{n-1}}$ and $\psi(x,c)$ is equivalent to a boolean combination of instances of $\varphi^{\text{opp}}$-formulas over $G_{\alpha^a_{n-1}}$, it follows by symmetry of $\varphi$ that $a\equiv^{\varphi^{\text{opp}}}_{G_{\alpha^a_{n-1}}} d_\gamma$ so $\psi(d_\gamma,c)$ holds as well.
%
%
%We claim that the collection of $d_\alpha$'s is a clique (of cardinality $\mu^+$), which is what we wanted.
%
%Let $\alpha_n^a<\alpha_1<\alpha_2<\mu^+$. Since the $\varphi$-types of $a$ and  $d_{\alpha_2}$  over $G_{\alpha_n^a}$ are defined by $\psi(x,c)$ we have $a\equiv^{\varphi}_{G_{\alpha_n^a}}d_{\alpha_2}$, and as $c\in G_{\alpha_0^a}$ and $\psi$ is a boolean combination of $\varphi^{\text{opp}}$-forumulas, we necessarily have that $\psi(d_{\alpha_2},c)$ holds. Consequently, $\varphi(d_{\alpha_1},d_{\alpha_2})$.
%%
%Let $p$ the $\varphi$-type $\psi(y,c)$ defines in $M$; set $d_0=a$, $d_1=b$ and choose elements in $M$ $d_n\vDash p|cd_{<n}$ for $2<n<\omega$. Since, for every $n<\omega$, $d_n\vDash p$, necessarily $\psi(d_n,c)$.
%
%We claim that $(d_n)_{n<\omega}$ form a complete subgraph of $M$. Indeed, for every $m<n<\omega$ with $n>2$, as $\psi(d_m,c)$ we conclude that $\varphi(x,d_m)\in p|cd_{<n}$ so $\varphi(d_n,d_m)$.
\end{proof}

Given Proposition \ref{P:stable} (and Proposition \ref{P:cannot embed infinite complete graph}), it is natural to ask under which conditions on the cardinality of $G$ does the  proposition hold. The following example shows that it fails for strong limit cardinals.

\begin{proposition}\label{P: stable no arbt large finite cliques}
Let $\lambda$ be a strong limit cardinal.\footnote{That means that $2^\mu<\lambda$ for any $\mu<\lambda$. Any such cardinal is a limit cardinal. An example is $\beth_{\omega}(\aleph_0)$.} There exists a non-stable graph with stable edge relation $G$ of cardinality $\lambda$ with $\chi(G)\geq \lambda$ for which we cannot embed arbitrary large finite cliques. In fact, it is triangle-free.
\end{proposition}
\begin{proof}
For any $\mu<\lambda$, let $G_\mu=\Sh_2(\beth_1(\mu)^+)$. By Lemma \ref{L:properties of shift2}, it is not stable but has a stable edge relation, and it is triangle-free. By Fact \ref{F:Sh-high chrom}, $\chi(G_\mu)\geq \mu^+$.

Let $G=\bigoplus_{\mu<\lambda}G_\mu$ be the direct sum of all of these graphs. Thus $\chi(G)\geq \lambda$ and $|G|=\lambda$. The graph $G$ is not stable but its edge relation is stable.  On the other hand, since each of the $G_\mu$'s is triangle-free,  we cannot embed arbitrary large finite cliques into $G$.
\end{proof}

\begin{remark} \label{R:even stable theory}
Using Remark \ref{R:Sh3}(2), we may replace $\Sh_2(\beth_1(\mu)^+)$ by $\Sh_3^{sym}(\beth_2(\mu)^+)$ and arrive at a stable graph with the prescribed properties.
\end{remark}




\section{The Chromatic Spectrum}\label{S:chromatic}
We rephrase the results of the previous section in a different manner.

For $T$ a theory of graphs and $\mu$ an infinite cardinal, let \[\ch(\mu)=\min\{ |G|: G\vDash T,\, \chi(G)\geq \mu\}.\]

We employ the convention that $\min\emptyset=\infty$. Note that if $\ch(\mu)=\infty$ then $\ch(\lambda)=\infty$ for all $\lambda\geq \mu$.

\begin{remark}
\begin{enumerate}
\item Since for any graph $G$, $\chi(G)\leq |G|$, $\ch(\mu)\geq \mu$.
\item Let $T=\mathrm{Th}(\Sh_k(\omega))$. By Fact \ref{F:Sh-high chrom}, $\ch(\lambda^+)\leq \beth_{k-1}(\lambda)^+$. On the other hand, if towards a contradiction we assume that $\ch(\lambda^+)<\beth_{k-1}(\lambda)^+$ then we can find $M\vDash T$ with $\chi(M)\geq \lambda^+$ and $|M|\leq \beth_{k-1}(\lambda)$. By L\"owenheim-Skolem we can can assume that $|M|=\beth_{k-1}(\lambda)$. By Lemma \ref{L:embed model of shift into symmetric}, we may embed $M$ into $\Sh_n^{sym}(\beth_{k-1}(\lambda))$, so by Fact \ref{F:Sh-high chrom} $\chi(M)\leq \lambda$, contradiction. We conclude that $\ch(\lambda^+)=\beth_{k-1}(\lambda)^+$.
%\item If $T$ is stable and $\{\chi(G): G\vDash T\}$ is bounded then $\chi(G)\leq \beth_2(\aleph_0)$ for all $G\vDash T$.
\end{enumerate}
\end{remark}

We can now rephrase the main result of \cite{1211} using this function:

\begin{proposition}
The following are equivalent for a stable theory of graphs  $T$:
\begin{enumerate}
\item There is some $G\vDash T$ and a natural number $k$ such that $G$ contains $\Sh_k(n)$ for all $n$.
\item $\ch(\beth_2(\aleph_0)^+)< \infty$.
\item For any cardinal $\mu$,  $\ch(\mu)< \infty$.
\item For any  cardinal $\mu$, $\ch(\mu)< \beth_\omega(\mu)$.
\end{enumerate}
\end{proposition}
\begin{proof}
$\neg (1)\implies \neg (2)$. By the main theorem of \cite[Corollary 6.2]{1211}, $\chi(G)\leq \beth_2(\aleph_0)$ for all $G\vDash T$, i.e. $\ch(\beth_2(\aleph_0)^+)=\infty$.

$(4)\implies (3)\implies (2)$. Easy.

$(1)\implies (4)$.  Let $\mu$ be an infinite cardinal. By compactness and (1) we can embed $\Sh_k(\beth_{k-1}(\mu)^+)$ in a large enough model of $T$. So by L\"owenheim-Skolem we can find a model $G$ of cardinality $\beth_{k-1}(\mu)^+$ with $\chi(G)\geq \chi\left(\Sh_k(\beth_{k-1}(\mu)^+\right)\geq \mu^+\geq \mu$ (using Fact \ref{F:Sh-high chrom}). Consequently, $\ch(\mu)\leq\beth_{k-1}(\mu)^+< \beth_\omega(\mu)$.
\end{proof}


Next, we phrase the results of the previous section for simple graphs and graphs with stable edge relation using $\ch$. 

\begin{proposition}\label{P:chi for simple and stable}
Let $T$ be a theory of graphs and assume that either $T$ is simple or the edge relation is stable. The following are equivalent:
\begin{enumerate}
\item $T$ proves the existence of arbitrary large finite cliques.
\item for any infinite cardinal $\mu$, $\ch(\mu)= \mu$.
\item for any infinite cardinal $\mu$, $\ch(\mu^+)= \mu^+$.
\item there exists an infinite cardinal $\mu$ for which $\ch(\mu^+)=\mu^+$.
\end{enumerate}
If $T$ is simple then they are also equivalent to:
\begin{enumerate}
\item[(5)] there exists an infinite cardinal $\mu$ with $\ch(\mu^+)\leq 2^\mu$.
\item[(6)] for any infinite cardinal $\mu$,  $\ch(\mu^+)\leq 2^{\mu}$.
\end{enumerate}

\end{proposition}
\begin{proof}
$(1)\implies (2)$.  By compactness and L\"owenheim-Skolem, there is a model $G$ of cardinality  $\mu$ which has an infinite clique of cardinality $\mu$. Thus $\chi(G)\geq \mu$, so $\ch(\mu)=\mu$.

 $(2)\implies (3)\implies (4)$ is easy.
 
$(4)\implies (1)$.  Assume that $(4)$ holds and let $\mu$ be an infinite cardinal for which $\ch(\mu^+)=\mu^+$. Thus there exists a model $G$ with $\chi(G)\geq \mu^+$ and $|G|= \mu^+$. By Proposition \ref{P:cannot embed infinite complete graph} for the simple case and Proposition \ref{P:stable} for stable edge relation case, $G$ contains arbitrary large finite cliques.

Assume that $T$ is simple.

$(4)\implies (5)$ is easy and $(5)\implies (1)$ uses Proposition \ref{P:cannot embed infinite complete graph} as above. 

$(3)\implies (6)\implies (5)$ is easy.
\end{proof}

%
%\begin{proposition} \label{prop:equivalence for simple}
%The following are equivalent for a simple theory of graphs $T$:
%\begin{enumerate}
%\item $T$ proves the existence of arbitrary large finite cliques.
%\item for any cardinal $\mu$, $\ch(\mu)=\mu$.
%\item for any cardinal $\mu$, $\ch(\mu^+)=\mu^+$.
%\item there exists a cardinal $\mu\geq 2^{\aleph_0}$ such that $\ch(\mu^+)=\mu^+$.
%%\item if $\mu\geq 2^{\aleph_0}$ then $\ch(\mu^+)>2^\mu$.
%\item[(5)] there exists a cardinal $\mu\geq 2^{\aleph_0}$ with $\ch(\mu^+)\leq 2^\mu$.
%\item[(6)] if $\mu\geq 2^{\aleph_0}$ then $\ch(\mu^+)\leq 2^{\mu}$.
%
%\end{enumerate}
%\end{proposition}
%\begin{proof}
%$(1)\implies (2)$.  By compactness and L\"owenheim-Skolem, there is a model $G$ of cardinality  $\mu$ which has an infinite clique of cardinality $\mu$. Thus $\chi(G)\geq \mu$, so $\ch(\mu)=\mu$.
%
%$(2)\implies (3)\implies (4)\implies (5)$  is easy.
%
%	 
%$(5)\implies (1)$ is Proposition \ref{P:cannot embed infinite complete graph}.
%
%$(3)\implies (6)\implies (5)$ is easy.
%\end{proof}
%
%
%Now for graphs with a stable edge relation.
%
%\begin{proposition} \label{prop:equivalence for stable edge}
%The following is equivalent for a theory of graphs $T$ with stable edge relation:
%
%\begin{enumerate}
%\item $T$ proves the existence of arbitrary large finite cliques.
%\item for any infinite cardinal $\mu$, $\ch(\mu)= \mu$.
%\item for any infinite cardinal $\mu$, $\ch(\mu^+)= \mu^+$.
%\item there exists an infinite cardinal $\mu$ for which $\ch(\mu^+)=\mu^+$.
%\end{enumerate}
%\end{proposition}
%\begin{proof}
%	$(1)\implies (2)$. By compactness there exists a model $G$ of cardinality  $\mu$ which has an infinite clique of cardinality $\mu$. Thus $\chi(G)\geq \mu$, so $\ch(\mu)=\mu$.
%	 
%	 $(2)\implies (3)\implies (4)$ is easy.
%	 
%	 $(4)\implies (1)$. Assume that $(4)$ holds and let $\mu$ be an infinite cardinal for which $\ch(\mu^+)=\mu^+$. Thus there exists a model $G$ with $\chi(G)\geq \mu^+$ and $|G|= \mu^+$. By Proposition \ref{P:stable}, $G$ contains an infinite clique (even a clique of cardinality $\mu^+$) giving (1). 
%\end{proof}
%
%We end we as an example of a theory of graphs for which $\ch(\mu)=\mu$ but is neither stable nor simple.
Is there an analog to Proposition \ref{P:chi for simple and stable} for general shift graphs? A reasonable suggestion is:
\begin{conjecture}
  Suppose that $T$ is stable. If for all cardinals $\mu$, $\ch(\mu^+) \leq \beth_{n-1}(\mu)^+$ then for some $m\leq n$, there is an embedding of $\Sh_{m}(\omega)$ in any $\omega$-saturated model of $T$.
\end{conjecture}
Note this is exactly Question \ref{Q:The question on Ch} without assuming GCH.

\appendix
\section{An Example by Hajnal and Komj\'{a}th} \label{A:Hajnal-Komjath example}

In this section we present an example due to Hajnal and Komj\'{a}th \cite[Theorem 4]{HK}. This is an example of a graph of size continuum whose chromatic number is $\aleph_1$, which does not contain all finite subgraphs of any shift graph $\Sh_n(\omega)$. They gave it as an example refuting Taylor's strong conjecture (which does hold outright for $\omega$-stable graphs with a close relative of it holding for stable theories in general by \cite{1196,1211}). The main goal here is to prove that this example has the independence property (IP) (thus is not stable) and furthermore that its theory is not simple.  


\begin{definition}\label{D:special}
  A graph $G=(V,E)$ is called \emph{special} if there exists partial order $\prec$ on $V$ satisfying that:
  \begin{enumerate}
    \item if $x\E y$ then either $x\prec y$ or $y\succ x$ and
    \item there is no circuit $C=\langle x_0,\dots,x_{n-1}\rangle$, $n\geq 3$, of the form \[x_0\prec x_1\prec \dots\prec x_{m-1}\prec x_m\succ x_{m+1}\succ \dots\succ x_{n-1}\succ x_0.\]
  \end{enumerate} 
\end{definition}


\begin{proposition}\label{P:special has no shift}\cite[Theorem 4]{HK}
Let $G=(V,E)$ be a special graph as witnessed by $\prec $. Then for all $n\geq 1$, $G$ does not contain all the finite subgraphs of $\Sh_n(\omega)$.
\end{proposition}
\begin{proof}
Assume towards a contradiction that $G$ contains all finite subgraphs of $\Sh_n(\omega)$ for some $n\geq 1$. If $n=1$ then it must contain a triangle so obviously contradicts Definition \ref{D:special}(1); so we assume that $n\geq 2$. Coloring pairs of $<$-increasing tuples, by Ramsey there is some integer $r$ such that if $f:\Sh_n(r)\to G$ is an embedding, there is $A\subseteq r$, $|A|=2n+1$ such that either $f(a_0,\dots,a_{n-1})\prec f(a_1,\dots,a_n)$ for all strictly increasing $n+1$-tuples from $A$ $(a_0,\dots,a_n)$ or $f(a_0,\dots,a_{n-1})\succ f(a_1,\dots,a_n)$ for all strictly increasing $n+1$-tuples from $A$, $(a_0,\dots,a_n)$. 

Assume the former occurs and that for simplicity $A=\{0,\dots,2n\}$. Then
\[(0,\dots,n-1)\prec (1,\dots,n-1,n+1)\prec\dots\prec (n-1,n+1,\dots, 2n-1)\prec (n+1,\dots,2n)\]
and
\[(0,\dots,n_1)\prec (1,\dots,n-1,n)\prec\dots\prec (n,n+1,\dots,2n-1)\prec (n+1,\dots,2n),\]
which is a contradiction.
\end{proof}

%
%\begin{fact}[L\'evy Collapse]\label{F:levy}
%Let $\mathbb{V}$ be a model of ZFC. If $\aleph_1<2^{\aleph_0}$ (in $\mathbb{V}$) then there exists a model $M\vDash $ZFC containing $\mathbb{V}$ satisfying:
%\begin{enumerate}
%\item $M$ and $\mathbb{V}$ have the same ordinals
%\item $\aleph_1=2^{\aleph_0}$ in $M$ and 
%\item for every set $X$ in $\mathbb{V}$, every function $f:\aleph_0\to X$ in $M$ belongs to $\mathbb{V}$.
%\end{enumerate}
%\end{fact}
%\begin{proof}
%By \cite[Lemma 15.21]{jech}, we may find such a forcing extension $M$ of $\mathbb{V}$ satisfying (2); in particular it satisfies (1) by \cite[Lemma 14.31(i)]{jech}. By the proof of \cite[Lemma 15.21]{jech}, the poset giving the forcing extension $M$ is $\aleph_1$-closed thus (2) follows by \cite[Fact 15.21]{halbeisen}.
%\end{proof}

%We reproduce the proof the their theorem here only to emphasize that indeed it has IP.
\begin{proposition}\label{P:Hajnal-Komjath Example}\cite[Theorem 4]{HK}
There exists a graph $G=(V,E)$, with $|V|=2^{\aleph_0}$ satisfying the following properties:
\begin{enumerate}
\item $G$ is special and in particular for every $n\geq 1$ it does not contain all finite subgraphs of $\Sh_n(\omega)$.
\item $\chi(G)=\aleph_1$.
%\item $G$ is not stable
\item $G$ has IP and in particular is not stable (in fact the edge relation has IP).
\item $G$ is not simple.
\end{enumerate}
\end{proposition}
\begin{proof}
Let $V=\{T_\alpha:\alpha<\aleph_1\}$ be a collection of disjoint sets with $|T_\alpha|=2^{\aleph_0}$ for each $\alpha<\aleph_1$. We define an edge relation on $V$ turning it to a graph satisfying our desired properties. 

To define the edge relation we define for $x\in T_\alpha$ and $\alpha<\aleph_1$, $G(x)=\{y\in T_{<\alpha}: x\E y\}$ by induction on $\alpha$.

If $\alpha=\beta+1$ we let $G(x)=\emptyset$ for every $x\in T_\alpha$ so assume that $\alpha$ is a limit ordinal and that $G(x)$ has already been defined for $x\in T_{<\alpha}=\bigcup_{\beta<\alpha} T_\alpha$.

For $\gamma<\alpha$ and $y\in T_{<\alpha}$, we say that \emph{$y$ is $\gamma$-covered} if there exists $\alpha_0<\dots<\alpha_m$, with $\alpha_0\leq \gamma$, and $x_i\in T_{\alpha_i}$ with $x_m=y$ such that $x_0\E x_1\dots \E x_m$.  Note that any $y\in T_\gamma$ is $\gamma$-covered as witnessed by the trivial path.

Let $\mathcal{W}_\alpha$ be the collection of all subsets $W\subseteq T_{<\alpha}$ satisfying that
\begin{list}{•}{}
\item $W=\{x_n: n<\omega\}$ is countable
\item $x_n\in T_{\alpha_n}$ and $\alpha_n<\alpha_m$ whenever $n<m<\omega$.
\item $\sup\{\alpha_n:n<\omega\}=\alpha$.
\item no $x_n$ is $\alpha_{n-1}$-covered for $0<n<\omega$.
\end{list}
Obviously, $|\mathcal{W}_\alpha|\leq 2^{\aleph_0}$; choose some enumeration $\mathcal{W}_\alpha=\{W_\gamma:\gamma<2^{\aleph_0}\}$ and $T_\alpha=\{t_\gamma:\gamma<2^{\aleph_0}\}$ and set $G(t_\gamma)=W_\gamma$. I.e for $y\in T_{<\alpha}$ and $x=t_\gamma\in T_\alpha$, $x\E y\iff y\in W_\gamma$. Let $G=(V,E)$.

We show that $G$ satisfies the properties listed in the statement.

We first show (1). Let $C=\langle x_0,\dots, x_{n-1}\rangle\subseteq V$, $n\geq 3$, be a circuit in $G$ with $x_i\in T_{\alpha_i}$ and
\[\alpha_0<\alpha_1<\dots<\alpha_{m-1}<\alpha_m>\alpha_{m+1}>\dots>\alpha_{n-1}>\alpha_0.\] for some $0<m<n-1$. If $\alpha_{m-1}= \alpha_{m+1}$, then as the elements of $C$ are distinct, $x_m$ would be connected to two vertices in $T_{\alpha_{m-1}}=T_{\alpha_{m+1}}$, contradicting our construction. So assume without loss of generality that $\alpha_{m-1}<\alpha_{m+1}$. Thus $x_{m-1},x_{m+1}\in G(x_m)$ and $x_{m+1}$ is $\alpha_{m-1}$-covered as witnessed by $\alpha_0<\alpha_{n-1}<\dots<\alpha_{m+1}$ and $\alpha_{m-1}\geq \alpha_0$; on the other hand since $\alpha_{m-1}<\alpha_{m+1}$ we get a contradiction.

We show (2). Let $c:G\to \aleph_0$ be a coloring. We say that a color $n<\omega$ is \emph{small} if there is a $\gamma_n<\aleph_1$ such that every point $x\in V$ with $c(x)=n$ is $\gamma_n$-covered; otherwise, call $n$ \emph{large}. For any small $n<\omega$ choose $\gamma_n$ minimal satisfying the above. Put $\gamma=\sup\{\gamma_n : n\text{ small}\}<\aleph_1$. We note that there exist large $n$; indeed take any $x\in T_{\gamma+1}$ and let $n<\omega$ be with $c(x)=n$. Since $x$ is not connected to any $y\in T_{<\gamma+1}$ it cannot be $\gamma$-covered. 

If $n<\aleph_0$ is large then for every $\alpha<\aleph_1$ there exists $x\in V$ with $c(x)=n$ which is not $\alpha$-covered. 

Let $\langle m_i<\aleph_0 :i<\omega\rangle$ be a sequence (possibly with repetitions), containing all large colors.

By definition, there exists $x_0\in T_{\alpha_0}$ with $c(x_0)=m_0$ which is not $\gamma$-covered (so necessarily $\gamma<\alpha_0$). We continue inductively and for any $n<\omega$ we find $x_n\in T_{\alpha_n}$ with $c(x_n)=m_n$ which is not $\alpha_{n-1}$-covered (so necessarily $\alpha_{n-1}<\alpha_n$). 

Let $\alpha=\sup\{\alpha_n:n<\omega\}$, it is necessarily a limit ordinal, and let $W=\{x_n :n<\omega\}\subseteq T_{<\alpha}$. By definition, there exists an element $x\in T_\alpha$ with $G(x)=W$. In particular $c(x)\neq m_n$ for all $n<\omega$, so $c(x)=k$ is small, i.e. it is $\gamma_k$-covered. But by the definition above, any $y\in G(x)=W$ is not $\gamma$-covered; so it cannot be that $x$ is $\gamma_k\leq \gamma$-covered.

To show that $\chi(G)=\aleph_1$ note that $c:V\to \aleph_1$ defined by $c(x)=\alpha_x$ for $x\in T_{\alpha_x}$ is a legal coloring.



%We now show (4), that $G$ is not stable. 
%
%Let $\mathbb{V}$ be a model of ZFC where $G$ is defined and let $\mathbb{V}\subseteq M$ be the model supplied by Fact \ref{F:levy}. Since $M$ and $\mathbb{V}$ have the same ordinals, and element $x\in T_\alpha$ is $\gamma$-covered in $\mathbb{V}$ if and only if it is $\gamma$-covered in $M$. Now, we need to re-produce the proof of $(2)$ inside $M$. By what we have just noted, being small or large does not in change when passing between $\mathbb{V}$ and $M$. In the proof we construct a countable subset of $V$; by the properties of $M$, this is still a countable subset in $\mathbb{V}$, so the proof can be concluded and we have thus showed that in $M$, $\chi(G)=\aleph_1=2^{\aleph_0}$.
%
%If $G$ were stable in $\mathbb{V}$ then it is also stable in $M$. By Lemma \ref{L:stable} and (3), since $|G|=\aleph_1$, we must have $\chi(G)\leq \aleph_0$; contradiction.

To show (3), choose for each $n<\omega$ some $x_n \in T_n$. Then, for every unbounded subsets $W\subseteq \omega$ there exists a unique $x\in T_\omega$ with $G(x)=\{x_n : n \in W\}$, giving IP.

Item (4) is a direct consequence of Proposition \ref{P:cannot embed infinite complete graph} (with $\mu=\aleph_0$).
\end{proof}
\begin{question}
Can one find such a counterexample which is stable? simple?
\end{question}


%
%\begin{remark}
%We remark here that we can also use Lemma \ref{L:stable} in order to show that the graph $G$ from Proposition \ref{P:Hajnal-Komjath Example} is not stable. 
%
%Let $\mathbb{V}$ be a model of ZFC where $G$ is defined and let $\mathbb{V}\subseteq M$ be a L\'evy collapse of $2^{\aleph_0}$ to $\aleph_1$. More specifically,  if $\aleph_1<2^{\aleph_0}$ (in $\mathbb{V}$) then there exists a model $M\vDash $ZFC (a forcing extension) containing $\mathbb{V}$ satisfying:
%\begin{enumerate}
%\item $M$ and $\mathbb{V}$ have the same ordinals
%\item $\aleph_1=2^{\aleph_0}$ in $M$ and 
%\item for every set $X$ in $\mathbb{V}$, every function $f:\aleph_0\to X$ in $M$ belongs to $\mathbb{V}$.
%\end{enumerate}
%To see this: by \cite[Lemma 15.21]{jech}, we may find such a forcing extension $M$ of $\mathbb{V}$ satisfying (2); in particular it satisfies (1) by \cite[Lemma 14.31(i)]{jech}. By the proof of \cite[Lemma 15.21]{jech}, the poset giving the forcing extension $M$ is $\aleph_1$-closed thus (2) follows by \cite[Fact 15.21]{halbeisen}.
%
%
%We use the notation of the proof of Proposition \ref{P:Hajnal-Komjath Example}. Let $M$ be as above. Since $M$ and $\mathbb{V}$ have the same ordinals, and element $x\in T_\alpha$ is $\gamma$-covered in $\mathbb{V}$ if and only if it is $\gamma$-covered in $M$. Now, we need to re-produce the proof of \ref{P:Hajnal-Komjath Example}(2) inside $M$. By what we have just noted, being small or large does not in change when passing between $\mathbb{V}$ and $M$. In the proof we construct a countable subsets of $V$; by the properties of $M$, these are still  countable subsets in $\mathbb{V}$, so the proof can be concluded and we have thus showed that in $M$, $\chi(G)=\aleph_1=2^{\aleph_0}$.
%
%If $G$ were stable in $\mathbb{V}$ then it is also stable in $M$. By Lemma \ref{L:stable}, since $|G|=\aleph_1$, we must have $\chi(G)\leq \aleph_0$; contradiction.
%\end{remark}
\bibliographystyle{alpha}
\bibliography{f2059}
\end{document}

%
%We define elements $a_{i,k,\eta}$ and sets $U_{i,\eta}\subseteq \mu$ such that $i\leq \mu$, $k\in U_{i,\eta}$ and $\eta\in (I^{\alpha_i})_\leq$.
%
%Assume we have defined elements with $i<j$, that are $(\alpha,I)$-indiscernible, and fix some $\eta\in (I^{\alpha_j})_\leq$. Let
%\[A_{j,\eta}=\dcl (\{a_{i,k,\nu}: i<j, k\in U_{i,\eta}, \nu\in (\eta^{\alpha_i})_\leq\}).\]
%Note that there are at most $\mu$ subsets $A\subseteq A_{j,\eta}$ with $|A|<\kappa$ and for all such $A$ there are at most $\mu$ non-principal ultrafilters on $A$. We view the set of pairs $U_{j,\eta}=\{(A,\mathcal{D})\}$ as a subset of $\mu$ and order them appropriately.
%
%For any other $\eta_0\in (I^{\alpha_j})_\leq$ the isomorphism $\pi=\pi_{\eta,\eta_0}:\eta\to \eta_0$ induces (by the induction hypothesis) a partial elementary map, which we also denote by $\pi$, $A_{j,\eta}\to A_{j,\eta_0}$, which is given by $a_{i,k,\nu}\mapsto a_{i,k,\pi(\nu)}$. Let $U_{j,\eta_0}=\{(\pi(A),\pi(\mathcal{D})),(A,\mathcal{D})\in U_{j,\eta}\}$.
%
%Order $(I^{\alpha_j})_\leq$ lexicographically. By induction on $k$, and by compactness, we may find for any $\eta\in (I^{\alpha_j})_\leq$,
%\[a_{j,k,\eta}\vDash p_{\mathcal{D}_k}|B_{j,k,\eta},\]
%where $\mathcal{D}_k$ is the ultrafilter on $A_k\subseteq A_{j,\eta}$ given by $k\in U_{j,\eta}$ and 
%\[B_{j,k,\eta}=\bigcup_{i<j,\nu\in (I^{\alpha_i})_{\leq}} A_{i,\nu}\cup \bigcup_{l\leq k,\, \nu<\eta}a_{j,l,\nu}\]
%
%To show $(\alpha,I)$-indiscernibility, assume it is known for elements $a_{i,l,\eta}$ with $(i,l)<(j,k)$. Let $\varphi(x_0,\dots,x_{n-1},\bar b)$ by a (quantifier-free) formula with $\bar b\subseteq \{a_{i,l,\eta}: (i,l)<(j,k),\,\eta\in (I^{\alpha_i})_\leq,\, l\in U_{i,\eta}\}$. Let $\eta_0>\dots>\eta_{n-1}\in (I^{\alpha_j})_\leq$ and let $\pi$ be a partial isomorphism of $(I,<)$ whose range contains all the relevant information. Note that by definition, \[\varphi(x_0,\dots,x_{n-1},\bar b)\in p_{\mathcal{D}_{j,\eta_0}}\otimes\dots\otimes p_{\mathcal{D}_{j,\eta_{n-1}}}\iff\]\[\varphi(x_0,\dots,x_{n-1},\pi(\bar b))\in p_{\mathcal{D}_{j,\pi(\eta_0)}}\otimes\dots\otimes p_{\mathcal{D}_{j,\pi(\eta_{n-1})}}.\]
%Indeed, if $\varphi(a_{j,k,\eta_0},\dots,a_{j,k,\eta_{n-1}},\bar b)$ then $\varphi(A,a_{j,k,\eta_2},\dots,a_{j,k,\eta_{n-1}},\bar b)\in \mathcal{D}$ for $(A,\mathcal{D})\in U_{j,\eta}$ corresponding to $k$. Since $\pi$ is elementary on $A\cup a_{j,k,\eta_2}\cup \dots\cup a_{j,k,\eta_{n-1}}\cup \bar b$, $\varphi(\pi(A),a_{j,k,\pi(\eta_2)},\dots,a_{j,k,\pi(\eta_{n-1})},\pi(\bar b))$.  The result follows.
%
%Now assume that $(I,<)$ is well ordered and let $A\subseteq \Rg(a)$ with $|A|<\kappa^+$. Let $\widehat \eta=\bigcup_{a_{i,\nu}\in A} \Rg(\nu)$ as a linearly ordered subset of $I$. (Don't we need $<\kappa$ so $\widehat \eta$ will be of length $<\kappa$?) Let $i<\mu$ be large enough indeed $(2^\kappa<\mu)$ such that $\alpha_i =\widehat \eta$ and $a_{j,\nu}\in A$ implies that $j\leq i$. This implies that $A\subseteq A_{i,\widehat \eta}$. Let $p\in S(A)$ and let $\tilde p$ be any extension of $p$ to $S(A_{i,\widehat \eta})$ (note that $A_{i,\widehat \eta}$ is a model).
